\chapter{文本分析}
\label{chap:nlp-text-analysis}

一份通知中的“4月18日”可以被准确地圈出来，但仅有位置还不能说明它是原定时间还是新时间。“校外来宾须报名”中的条件作用于谁，也不能由“来宾”和“报名”两个关键词独自回答。第~\ref{chap:nlp-introduction}章已经区分任务输出与判断依据；现在需要把这些区分落实为能够恢复、检查和传给后续任务的分析结果。

文本分析从五个方面回答这个问题。词法确定语言单位及局部属性，句法约束单位怎样组合，语义选择词义并表达谓词、论元或形式结构，篇章连接跨句对象与内容，语用则利用交流背景解释话语的作用。它们是分析视角，不是一条必须顺序执行的流水线：上文出现的身份可能改变当前词义，交流背景也可能排除一种句法解释。下游系统可以使用其中某些显式结果，也可以直接从文本预测自己的输出。

本章所有短句、概率、得分、数据库和动作轨迹均为教学构造，不是模型实测。沿用从零开始、左闭右开的原文字符跨度$[s,e)$；字符索引按Unicode字符计数，不是UTF-8字节偏移。数学计算中的位置可从1开始，但每次恢复结果都回到原文偏移。语言词语、模型词元和内部向量分别命名，不能把向量存在当作结构已经形成。

\section{词法分析与局部标注}
\label{sec:nlp-analysis-lexical}

“机器学习读书会”可以作为一个活动名称，但词性标注可能需要“机器学习／读书会”，模型输入又可能把每个汉字作为词元。\term{词法分析}（lexical analysis）处理词的边界、形式和局部语法属性；它必须先确定采用什么语言单位，再给这些单位赋属性。一个词被拆成两个词后，词级标签即使都合理，也未必能与原标注逐位置比较。

图~\ref{fig:nlp-analysis-alignment}把位置作为各层共同的参照。图中的模型切分只是教学设定，不对应某个实际词表。分词和组块可以有不同边界，模型词元也不必服从其中任何一种边界。

\begin{figure}[htbp]
  \centering
  \begin{tikzpicture}[
  x=1mm,y=1mm,
  every node/.append style={font=\figurefont\footnotesize,align=center},
  unit/.style={draw=FigureOutput,fill=FigureOutput!10,line width=.85pt,
    minimum height=9mm,inner sep=1mm},
  mapping/.style={->,draw=NNLine,line width=.85pt}
]
  \node[anchor=east] at (21,17) {模型词元};
  \node[anchor=east] at (21,0) {原文字符};
  \node[anchor=east] at (21,-17) {语言词};
  \node[anchor=east] at (21,-31) {BMES};
  \node[anchor=east] at (21,-42) {词级组块};
  \node[unit,draw=FigureModel,fill=FigureModel!10,minimum width=51mm]
    at (50,17) {研究生 \quad $[0,3)$};
  \node[unit,draw=FigureModel,fill=FigureModel!10,minimum width=17mm]
    at (84,17) {命};
  \foreach \x/\char/\offset in {33/研/0,50/究/1,67/生/2,84/命/3} {
    \node[unit,draw=FigureInput,fill=FigureInput!10,minimum width=17mm]
      at (\x,0) {\char};
    \node[font=\figurefont\scriptsize,anchor=south] at (\x,4.8) {\offset};
    \draw[mapping] (\x,12.5) -- (\x,9);
    \draw[mapping] (\x,-12.5) -- (\x,-4.5);
  }
  \node[unit,minimum width=34mm] at (41.5,-17) {研究 \quad $[0,2)$};
  \node[unit,minimum width=34mm] at (75.5,-17) {生命 \quad $[2,4)$};
  \foreach \x/\tag in {33/B,50/E,67/B,84/E}
    \node at (\x,-31) {\tag};
  \node[unit,minimum width=34mm] at (41.5,-42) {B-VP};
  \node[unit,minimum width=34mm] at (75.5,-42) {B-NP};
  \node[anchor=north,align=left] at (58.5,-50)
    {组块：研究/VP，生命/NP\\箭头：单位覆盖的原文字符};
\end{tikzpicture}

  \caption{原文“研究生命”的字符、语言词、模型词元与标签对应。横向位置均指原文字符；箭头只表示覆盖位置，不表示预测依赖。语言词采用“研究／生命”，模型词元故意采用“研究生／命”。本例把“研究”视为动词、“生命”视为名词，分别形成VP与NP组块；词元标签必须经过位置映射才能恢复这些语言单位。}
  \label{fig:nlp-analysis-alignment}
\end{figure}

\subsection{分词与边界确定}
\label{sec:nlp-analysis-segmentation}

\term{分词}（word segmentation）按照选定规范把字符序列分成词语。中文没有用空格普遍标出词界，“研究生命”既包含词典中的“研究生”，也包含“研究”和“生命”；局部最长匹配未必对应当前语境。领域术语“机器学习”、时间串“14:00”和补充例子中的人名“小王”，还分别涉及术语是否合并、数字标点怎样处理及专名词典的覆盖。规范必须先于训练与评价固定下来。

语言学分词不等于模型的子词切分。BPE按合并规则，WordPiece按词表及其切分规则，Unigram分词模型按候选词元序列概率组织输入；这些目标不会自动生成所需语言词界，输入表示的背景可回看第~\ref{chap:transformer}章。固定“研究生命”时，字符输入为四个位置，语言词可为“研究／生命”，模型词元可为“研究生／命”。后者跨过了语言词界，不能把一个词元的类别直接复制成一个完整词的类别。

评价时将预测词和参考词都恢复为原文跨度集合。预测“研究生／命”得到$\{[0,3),[3,4)\}$，参考“研究／生命”为$\{[0,2),[2,4)\}$，严格词匹配数为0。即使某个下游实体仍能由$[0,4)$正确恢复，这也不能把分词错误改记为正确；应另报下游的端到端结果。

\subsubsection{词典最大匹配与候选切分图}
\label{sec:nlp-analysis-dictionary}

设词典包含“研究”“研究生”“生命”“命”，并允许单字兜底。正向最大匹配从左边当前位置出发，尝试不超过词典最大词长的片段，取最长的已收录词，然后移动到其右端。因此“研究生命”得到“研究生／命”。逆向最大匹配从右端寻找最长词，先选“生命”，再选“研究”，得到另一种切分。两种方法都只保留当前局部选择；所谓双向比较还需要额外规定冲突时按词数、单字数或其他证据选择，不能由“最大匹配”本身推出答案。

要同时比较整句候选，可把字符间的边界$0,\ldots,n$作为节点。若$x[s:e]$在词典中，就加入边$s\to e$，边载荷为该词及其代价$c(s,e)$。所有边都向右，因此构成有向无环图。令$d(e)$为从0到$e$的最小总代价，递推为
\begin{equation}
  d(0)=0,\qquad
  d(e)=\min_{s:\,s\to e}\{d(s)+c(s,e)\}.
  \label{eq:nlp-analysis-word-path}
\end{equation}
其余初值为$+\infty$，每次保存最优前驱$s$。到达$n$后逆向读取前驱，把各条边的$[s,e)$反转为正序，就是词跨度。若所有到达$n$的候选都不可达，应返回无合法切分，或明确启用单字兜底，而不是悄悄跳过字符。

在本例中，设“研究”“生命”“研究生”“命”的代价依次为$1,1,0.7,2$，单字兜底代价均为3。于是$d(2)=1$，$d(3)=0.7$；到达4时，经“生命”的代价是$1+1=2$，经“命”的代价是$0.7+2=2.7$，其余路径更大。回溯$4\leftarrow2\leftarrow0$得到“研究／生命”。若代价取负对数词概率，就成为最大乘积概率路径；这些概率须从训练语料估计，不能将本例任意给定的代价称为训练结果。词典缺少“生命”会直接删掉正确候选，而错误长词会使贪心提前走错，两种故障应分别修复候选覆盖和路径评分。

\subsubsection{HMM边界标注}
\label{sec:nlp-analysis-hmm-boundary}

词典不能枚举所有新词，字符在词内的位置却可以共享。\term{BMES标注}分别用B表示词首、M表示词中、E表示词尾、S表示单字词；“研究／生命”转成B、E、B、E，“研究生／命”转成B、M、E、S。这样，一个未见过的词仍可由已经见过的边界模式组成。

隐马尔可夫模型（HMM）把边界标签$y_i$当作隐藏状态，字符$x_i$当作观测。其条件独立性和前向推断见第~\ref{sec:pgm-dynamic-discrete}节；这里的参数来源是已分词文本转换出的BMES序列。训练时统计句首状态次数$N_{\mathrm{start}}(a)$、相邻状态次数$N(a,b)$以及状态发射字符次数$N(a,c)$。例如用加$\alpha$平滑，在允许后继集合$A(a)$上估计
\[
  p(b\mid a)=
  \frac{N(a,b)+\alpha}
       {\sum_{b'\in A(a)}N(a,b')+\alpha|A(a)|}.
\]
发射概率也按状态分别归一化，并在字符词表中显式保留未知字符符号。句末可作为E、S的额外后继$\mathrm{END}$参与计数；非法转移始终是0，不参与平滑分母。

考虑两条训练记录“研究／生命”和“研究生／命”。它们贡献B到E两次、B到M一次，因此不加平滑时$p(E\mid B)=2/3$、$p(M\mid B)=1/3$。在第一条记录中，“研”和“生”贡献B发射；第二条中“研”贡献B、“生”贡献E。可见同一个字符能处在不同词内位置，观测词频与边界转移提供的是不同信息。稀疏语料若把未见“起”字的发射概率记成0，“研究起源”的所有路径可能都被消掉；未知字符平滑处理的是这个问题，并不能补出未知词的意义。

固定待分析的四个字符，设一个已平滑模型给两条候选的连乘因子如下：
\[
\begin{aligned}
 p(x,\mathrm{BEBE})&=
 .8\cdot.5\cdot.6\cdot.3\cdot.6\cdot.3
 \cdot.6\cdot.4\cdot.2=0.00062208,\\
 p(x,\mathrm{BMES})&=
 .8\cdot.5\cdot.4\cdot.5\cdot.9\cdot.2
 \cdot.2\cdot.2\cdot.8=0.0004608.
\end{aligned}
\]
因子依次为起始概率、字符发射、转移与发射的交替，以及句末转移。这些是另一组教学参数，不是前面两条记录的估计值。比较这两条路径时BEBE胜出；要确认整句最优，仍须让Viterbi覆盖所有合法路径，而不能只比较两个手选候选。取对数后，下一小小节的递推可直接使用发射对数概率和转移对数概率，最终恢复$[0,2)$、$[2,4)$。概率估计是否合理、搜索是否正确、BMES是否合法，是三项独立检查。

\subsubsection{BMES--CRF边界标注与Viterbi解码}
\label{sec:nlp-analysis-bmes-crf}

HMM的发射分布只直接依赖当前状态。若要同时使用左右邻字、词典命中、字符类型或全句表示，可以直接为“这个位置取某个标签”评分，再用标签间的得分连接整句。\term{线性链条件随机场}（linear-chain conditional random field, CRF）的通用归一化与梯度见第~\ref{sec:pgm-energy-models-crf-definition}节。这里需要补齐的是，语言边界规范怎样进入它的输出空间。

设句长为$n$，标签集合$Y=\{\mathrm{B,M,E,S}\}$。位置得分$e_i(a)$可以是特征线性组合$\bm{w}_a^\top\bm{f}(x,i)$，也可以由神经编码器的上下文向量$\bm{h}_i$经仿射变换得到。可学习转移矩阵为$\bm{A}$；起止得分分别为$u(a)$、$v(a)$。合法性另用固定掩码规定：
\[
\begin{array}{c|c}
 \text{位置或前驱}&\text{允许标签或后继}\\ \hline
 \text{句首}&\mathrm{B,S}\\
 \mathrm{B}&\mathrm{M,E}\\
 \mathrm{M}&\mathrm{M,E}\\
 \mathrm{E}&\mathrm{B,S}\\
 \mathrm{S}&\mathrm{B,S}\\
 \text{句末}&\mathrm{E,S}
\end{array}
\]
把非法转移及非法起止得分设为$-\infty$，合法项保持原得分，记掩码后的参数为$\widetilde{\bm{A}},\tilde u,\tilde v$。对合法标签序列集合$\mathcal{Y}(x)$，整句得分和条件概率是
\begin{align}
 S(x,y)&=\tilde u(y_1)+\sum_{i=1}^{n}e_i(y_i)
       +\sum_{i=2}^{n}\widetilde{\bm{A}}_{y_{i-1},y_i}
       +\tilde v(y_n),\label{eq:nlp-analysis-crf-score}\\
 p(y\mid x)&=\frac{\exp S(x,y)}
   {\sum_{y'\in\mathcal{Y}(x)}\exp S(x,y')}.\label{eq:nlp-analysis-crf-probability}
\end{align}
训练样本是字符序列和由参考词界转换的标签序列$y^\star$，损失为$\log Z(x)-S(x,y^\star)$。训练时的$Z$也必须只累加合法路径；若参考标签违反规范，应修正标注或转换程序，不能让它进入这个目标。梯度提高参考序列的特征或神经得分，并扣去模型分布下相应的期望贡献。它更新的是评分参数，固定合法性掩码不需要学习。

归一化可用前向递推实现。令$\alpha_1(a)=\tilde u(a)+e_1(a)$，则
\[
\begin{aligned}
 \alpha_i(b)&=e_i(b)+
 \operatorname{logsumexp}_{a\in Y}
       [\alpha_{i-1}(a)+\widetilde{\bm{A}}_{a,b}],\\
 \log Z&=\operatorname{logsumexp}_{a\in Y}
       [\alpha_n(a)+\tilde v(a)].
\end{aligned}
\]
全部候选为$-\infty$时，累加结果仍是$-\infty$。实现中不能直接计算$\exp$后相加，应使用稳定的log-sum-exp。配分函数是在比较所有合法序列，不是只比较Viterbi的赢家。

解码需要最高分序列。令$\delta_i(b)$为结束于位置$i$、标签$b$的最高前缀分数。具有同一末标签的两个前缀，未来新增得分相同，因此只保留其中较大者不会丢掉全局最优。由此得到
\begin{align}
 \delta_1(b)&=\tilde u(b)+e_1(b),\nonumber\\
 \delta_i(b)&=e_i(b)+\max_{a\in Y}
   [\delta_{i-1}(a)+\widetilde{\bm{A}}_{a,b}],\label{eq:nlp-analysis-viterbi}\\
 \psi_i(b)&=\argmax_{a\in Y}
   [\delta_{i-1}(a)+\widetilde{\bm{A}}_{a,b}].\nonumber
\end{align}
先选$y_n=\argmax_b[\delta_n(b)+\tilde v(b)]$，再从$i=n$到2反复令$y_{i-1}=\psi_i(y_i)$。平局时固定标签顺序，保证结果可复现；平局本身意味着存在多个同分候选，而非一个已经确定的语言事实。时间为$O(n|Y|^2)$，保存回溯指针需要$O(n|Y|)$空间，编码器成本另计。

对“研究生命”，令所有合法转移、起止得分为0，使用表~\ref{tab:nlp-analysis-viterbi}的位置得分。逐位置选择最大值会得到SEBE，分数为15，但首个S后不能接E。强行把最高的四个局部标签拼起来，根本不在允许的输出空间中。

\begin{table}[htbp]
  \centering
  \begin{tabular}{@{}c rrrr rrrr@{}}
    \toprule
    &\multicolumn{4}{c}{位置得分$e_i$}
    &\multicolumn{4}{c}{最优前缀$\delta_i$}\\
    字&B&M&E&S&B&M&E&S\\
    \midrule
    研&3&0&0&4&3&$-\infty$&$-\infty$&4\\
    究&0&1&4&0&4&4&7&4\\
    生&3&1&2&0&10&5&6&7\\
    命&0&0&4&1&7&10&14&8\\
    \bottomrule
  \end{tabular}
  \caption{BMES的教学得分与Viterbi图表。起止和转移均采用正文掩码，合法项为0。右半表按末标签分别保留最优前缀，因此句末的M虽然有10分，仍不允许结束。}
  \label{tab:nlp-analysis-viterbi}
\end{table}

例如$\delta_2(\mathrm E)=4+\delta_1(\mathrm B)=7$，而$\delta_3(\mathrm B)=3+\max\{\delta_2(\mathrm E),\delta_2(\mathrm S)\}=10$。句末只比较E的14和S的8，选择E。其前驱依次为位置3的B、位置2的E、位置1的B，得到BEBE，分数$3+4+3+4=14$。恢复时，B开启跨度，M延长跨度，E闭合跨度，S独自产生长度1的词；按原文偏移输出“研究”$[0,2)$和“生命”$[2,4)$。训练允许的结构、推断使用的结构和恢复程序接受的结构由同一张转移表约束。

这个例子也说明三种常见错误的区别。未知词不是禁止输出的词典缺项，只要字符表示和边界组合足够，仍能恢复新词；领域术语若按训练规范被切得很细，则需要新标注或领域特征，不能只改回溯程序；原文经过规范化而偏移未保存时，即使BEBE正确也会返回错误位置。长度1的输入只能取S，空输入可按接口约定返回空词表；如果非空输入的最终分数全为$-\infty$，应报告无合法路径。用有限的大负数近似掩码时还要检查数值范围，避免非法路径因溢出或精度问题重新胜出。

\subsection{词形还原与形态属性}
\label{sec:nlp-analysis-morphology}

得到词界后，同一个词还可能具有多个表面形式。“studies”“studied”通常可还原为“study”，但“studies”在句中可能是复数名词，也可能是动词第三人称单数形式。\term{基本形式}（lemma）是词典记录或标注规范选定的代表形式；\term{形态属性}（morphological features）记录数、时态、人称等语法变化。基本形式不是模型词元，也不是简单截掉词尾后留下的字符串：把“studies”截成“studi”只是词干提取可能产生的结果。

这些任务依赖语言。中文“昨天讨论”和“明天讨论”不靠动词屈折区分过去与将来，不能强行给“讨论”恢复一个英语式过去时词尾。语码转换，即一次表达中切换语言，还要求识别每个片段的语言及适用规则。以下三个方法都输出基本形式，但形态属性必须有独立标注与恢复步骤；评价分别检查基本形式精确匹配、逐项属性以及整个属性组合。

\subsubsection{词典与有限状态形态分析}
\label{sec:nlp-analysis-fst}

只查表面词形会错过许多规则变化。\term{有限状态转导器}（finite-state transducer, FST）在一条状态路径上同时读取一种符号串、输出另一种符号串，因而可以把“基本形式＋属性”与表面形式联系起来。词典规定词干和允许的词缀，规则规定拼写变化；沿反向关系查询表面词形，就获得可能的分析，而不是唯一答案。

构造一个只覆盖英语辅音加y词尾的微型规则：词干\code{stud}后接词尾\code{y}，当属性为复数名词或动词第三人称单数现在时，输出词尾\code{ies}。同一输入\code{studies}因此具有两条接受路径：
\[
\begin{aligned}
 \code{studies}&\leftrightarrow
   (\code{study},\mathrm{NOUN},\mathrm{Number=Plur}),\\
 \code{studies}&\leftrightarrow
   (\code{study},\mathrm{VERB},\mathrm{Tense=Pres},
                       \mathrm{Person=3},\mathrm{Number=Sing}).
\end{aligned}
\]
实现可以把词干路径与后缀路径连接，再与拼写规则转导器复合；本例也可用等价的小规则表手工展开。接受状态只说明路径允许这一映射，不能说明它符合当前句子。

在“She studies”中，为两个候选提取“左邻词为代词”“主语第三人称单数”等特征，设分类器给名词候选$-1$、动词候选2，softmax后动词概率约为0.953；“two studies”中若得分反为2与$-1$，则选复数名词。分类器由语境和正确分析对训练，FST规则由语言资源提供，这两种知识来源不能混写成同一次训练。若删除词干\code{stud}，是词典漏候选；若不允许y到ies，是拼写规则漏候选。未见形式可返回未知或交给下一种方法，但不得把互斥的单数、复数属性同时附在同一分析上。

\subsubsection{词形编辑脚本预测}
\label{sec:nlp-analysis-edit-lemma}

大量词语共享相同变化时，可以预测“怎样修改词形”，不必给每个完整词形单独建类别。\term{编辑脚本}（edit script）是一段可执行的字符修改操作。将训练词形与基本形式对齐，保留共同部分，再把删除、替换和大小写调整组织成脚本；对齐平局需有固定规则，否则同一训练对会产生不一致监督。

对\code{studies}与\code{study}，一个相对词尾的脚本是“保留前$|w|-3$个字符，删去末尾\code{ies}，附加\code{y}”。用语境表示$\bm h$预测脚本$r$，得分$g_r=\bm w_r^\top\bm h+b_r$，以正确脚本的交叉熵训练。设该脚本、恒等脚本、删除末尾s三个得分为$2,0,-1$，归一化后概率约为$0.844,0.114,0.042$；选择第一项并真正执行，才得到\code{study}。

对未见过的\code{retries}，同一脚本保留\code{retr}并补y，恢复\code{retry}，说明类别可跨词共享。它却会把\code{ties}错误地变成\code{ty}；正确基本形式是\code{tie}，需要另一脚本及语境选择。执行前应检查长度和后缀前提，执行后再查目标词典或合法字符条件。脚本类别预测正确、脚本执行无越界、结果符合目标语言是不同检查；词典未收录的新词也不能仅因查无结果就一律判错。

\subsubsection{字符级序列到序列还原}
\label{sec:nlp-analysis-seq2seq-lemma}

编辑脚本集合若没有所需变换，可以直接生成基本形式的字符序列。编码器读取表面字符及允许的词语上下文，解码器在已输出字符条件下预测下一个字符，直到结束符。对训练对$(w,\ell^\star)$，包括结束符的条件负对数似然为
\[
  -\sum_{j=1}^{|\ell^\star|+1}
  \log p(\ell_j^\star\mid w,\text{上下文},\ell^\star_{<j}).
\]
这是前卷编码器--解码器在字符层面的使用，不要求输入输出长度相同。保留字符可以由普通生成学习，也可显式混合词表生成与输入位置复制；复制分布中同字符的多个位置必须相加。

仍以\code{retries}为输入，设正确路径依次输出r、e、t、r、y、结束符的条件概率为$.9,.9,.8,.8,.7,.95$，路径概率为$0.344736$。假设复制输入并直接输出\code{retries}的候选路径概率为0.12，束搜索保留两者并在完成后选择前者，恢复\code{retry}。这里给的是条件路径概率，不能将各位置在不同历史下的最大概率任意相乘。长度上限到了却没有结束符，应标记截断，不能把残缺输出视为正常基本形式。

另设属性预测头，为时态、人称、数和词性各预测一个包含“不适用”的类别；或采用多标签输出后进行合法组合检查。对于“She retries”，输出基本形式\code{retry}的同时，可恢复$\mathrm{VERB,Pres,3,Sing}$；对于“two retries”则是$\mathrm{NOUN,Plur}$。相同基本形式并不决定这些属性。字符生成能组合未见串，也会拼出不存在的形式；需检查遗漏、重复、拼写和上下文，而不能用路径概率代替语言分析正确性。

\subsection{词性标注与短语组块}
\label{sec:nlp-analysis-pos-chunk}

\term{词性标注}（part-of-speech tagging, POS tagging）判定一个词在当前语境中的语法类别；“研究问题”中的“研究”通常用作动词，“这项研究”中则用作名词。\term{短语组块}（chunking）只标出浅层短语范围，例如一个连续名词短语，不递归恢复全部嵌套结构。两者分别输出类别序列和带类别的区间，后者不能用逐词准确率代替完整组块F1。

给定“校外／来宾／须／报名”的词界，逐词线性分类器可以用邻词、词形和词典特征评分，通过softmax选择标签；一般分类损失见第~\ref{chap:classification}章。若“研究”的名词、动词得分为$0.4,0.6$，独立选择只得到动词，不会检查其余标签是否形成合理组合。下面的方法分别增加状态依赖、全句表示或联合边界输出。若起点是原始文本，还必须在原文跨度上比较词性和组块，不能假定错误切分后的第$i$个词仍对应参考的第$i$个词。

\subsubsection{HMM词性标注}
\label{sec:nlp-analysis-hmm-pos}

把第~\ref{sec:nlp-analysis-hmm-boundary}节的字符观测换成词语观测、BMES状态换成词性，就得到HMM词性标注。由已标注词序列统计$p(t_i\mid t_{i-1})$、$p(w_i\mid t_i)$和起止概率，Viterbi对完整词性路径求最大值。词性转移并没有BMES那样普适的禁止表；除了任务规定的起止条件，低频组合不应自动视为语法不可能。

考虑给定分词的“研究／问题”，只比较动词--名词（V、N）和名词--名词（N、N）。设起始概率$p(V)=.6,p(N)=.4$，观测概率$p(\text{研究}\mid V)=.3,p(\text{研究}\mid N)=.4$，转移$p(N\mid V)=.8,p(N\mid N)=.2$，以及$p(\text{问题}\mid N)=.5$；本例两条路径的句末因子相同而省略。两条路径得分分别为$.6\cdot.3\cdot.8\cdot.5=.072$与$.4\cdot.4\cdot.2\cdot.5=.016$。虽然“研究”的名词观测概率更高，完整序列仍选择V、N。恢复结果为“研究”$[0,2)$：V，“问题”$[2,4)$：N。

遇到未见词，可把低频词在训练时映射成含数字、后缀或大小写特征的未知词类别，再估计该类别的状态发射分布。不能只在测试时替换却沿用一个从未训练的观测值。以上结果以词界已给定为前提；原文切成“研究问／题”时，再高的标签准确性也不修复词界。

\subsubsection{CRF与BiLSTM-CRF局部标注}
\label{sec:nlp-analysis-crf-chunk}

特征CRF用左右邻词和词形特征产生位置得分，BiLSTM-CRF则用双向循环网络产生上下文表示，再交给同一个CRF目标和解码器。双向表示可使用当前句子的后文；它是离线句子分析的可用信息，不应在只能看到输入前缀的任务中原样使用。编码器改变得分的来源，标签集合和合法输出并未因此改变。

对浅层名词组块，采用BIO标签：B-NP开始组块，I-NP延续同类组块，O表示外部。句首禁止I-NP，I-NP只能接在B-NP或I-NP之后；多类组块还要求延续标签的类型相同。以“校外／来宾／须／报名”为例，前三种标签得分分别设为
\[
\begin{array}{c|rrr}
 &\mathrm{B\!-\!NP}&\mathrm{I\!-\!NP}&\mathrm O\\ \hline
 \text{校外}&3&4&0\\
 \text{来宾}&0&4&1\\
 \text{须}&0&0&3\\
 \text{报名}&0&0&3
\end{array}
\]
合法转移为0时，独立最大值为I-NP、I-NP、O、O，首标签非法；约束Viterbi恢复B-NP、I-NP、O、O，总分13。读到第一个O时闭合名词组块，映射词的起止字符位置，得到“校外来宾”$[0,4)$。

同样地，“研究”在“研究问题”和“这项研究”中可获得不同的上下文得分。特征CRF依赖人工指定的邻词特征，独立神经分类器虽能看全句却不约束输出序列，BiLSTM-CRF同时提供上下文评分和结构解码。三者应在相同词界、标签规范和评价样本下比较。BIOES可以显式区分组块尾部和单词组块，但需要相应修改起止和转移表，不能把BIO掩码原封不动复用。

\subsubsection{字符级分词--词性联合标注}
\label{sec:nlp-analysis-joint-pos}

流水线一旦把“研究生”固定成词，后面的词性标注器便无法恢复“研究／生命”。联合标注把边界与词性组合为一个标签，例如B-V、E-V、B-N、E-N，使词性证据能够参与词界选择。训练时，将每个已标注词展开为字符标签；“研究/V，生命/N”产生B-V、E-V、B-N、E-N，“研究生/N，命/N”则产生B-N、M-N、E-N、S-N。

沿用CRF计算，但B-$t$和M-$t$只允许转移到M-$t$或E-$t$，保证同一词内词性一致；E-$t$、S-$t$才允许转到另一词性的B或S。假设两条候选的完整得分为8和6，其余路径低于6，约束解码选择前者。恢复程序在E-V处闭合“研究”$[0,2)$并附V，在E-N处闭合“生命”$[2,4)$并附N。B-V、E-N这样的混合词内类别即使局部分数高，也被掩码排除。

若共有$k$种词性，联合标签最多为$4k$，直接稠密递推的代价随标签数平方增长，合法转移稀疏化可以减少无用比较。联合训练需要同时有词界和词性监督，不能从只有分词的数据自动得到正确词性。恢复出的两个词仍没有说明它们是否组成一个短语组块；组块必须按自己的边界与类别规范另行预测。联合输出减少了一处硬流水线决策，却没有保证所有局部分析都正确。

\section{句法分析}
\label{sec:nlp-analysis-syntax}

局部标签没有给出短语如何嵌套，也没有说明哪个词支配哪个词。\term{成分结构}（constituency structure）用树表示相邻词如何组成更大的短语；\term{依存结构}（dependency structure）用有向边表示中心词与依存词的关系。前者的内部节点是短语，后者的基本节点是词，不能因两者都画成树就交换其节点含义。

图~\ref{fig:nlp-analysis-syntax}固定“小王／阅读／通知”的分词。成分树将“阅读通知”组成动词短语，再与“小王”组成句子；依存树以“阅读”为中心，连接主语“小王”和宾语“通知”。两种分析可以通过额外的中心词规则联系，但通常不是无损的一一互换，更不是先后执行的两个阶段。

\begin{figure*}[htbp]
  \centering
  \begin{tikzpicture}[
  x=1mm,y=1mm,
  every node/.append style={font=\figurefont\footnotesize,align=center},
  category/.style={draw=FigureModel,fill=FigureModel!10,
    rounded corners=1mm,inner sep=1.5mm},
  word/.style={draw=FigureInput,fill=FigureInput!10,
    minimum width=16mm,minimum height=8mm,inner sep=1mm},
  branch/.style={draw=NNLine,line width=.85pt},
  dep/.style={->,draw=NNWire,line width=1.1pt}
]
  \node[font=\figurefont\footnotesize\bfseries] at (40,59) {(a) 成分：短语包含关系};
  \node[category] (s) at (39,46) {S};
  \node[category] (np1) at (13,30) {NP};
  \node[category] (vp) at (53,30) {VP};
  \node[category] (v) at (41,15) {V};
  \node[category] (np2) at (68,15) {NP};
  \node[word] (a) at (13,0) {小王};
  \node[word] (b) at (41,0) {阅读};
  \node[word] (c) at (68,0) {通知};
  \foreach \parent/\child in {s/np1,s/vp,np1/a,vp/v,vp/np2,v/b,np2/c}
    \draw[branch] (\parent.south) -- (\child.north);
  \node at (13,-10) {$[0,2)$};
  \node at (41,-10) {$[2,4)$};
  \node at (68,-10) {$[4,6)$};

  \node[font=\figurefont\footnotesize\bfseries] at (127,59) {(b) 依存：词间有向关系};
  \node[category] (root) at (127,43) {0：人工根};
  \node[word] (d) at (102,0) {小王};
  \node[word] (e) at (127,0) {阅读};
  \node[word] (f) at (152,0) {通知};
  \draw[dep] (root.south) -- node[right,fill=white,inner sep=1mm] {root} (e.north);
  \draw[dep] (e.north west) to[out=145,in=45]
    node[above left,inner sep=1mm] {nsubj} (d.north);
  \draw[dep] (e.north east) to[out=35,in=135]
    node[above right,inner sep=1mm] {obj} (f.north);
  \node at (102,-10) {$1:[0,2)$};
  \node at (127,-10) {$2:[2,4)$};
  \node at (152,-10) {$3:[4,6)$};
  \draw[NNLine,dashed,line width=.5pt] (84,-14) -- (84,52);
\end{tikzpicture}

  \caption{同一句话的成分树与基本依存树。左图的NP、VP、S分别表示名词短语、动词短语和句子；右图的边由中心词指向依存词，0是人工根，nsubj和obj分别标记主语与宾语关系。这里固定分词和一种教学分析，不把两种表示画成处理流水线。}
  \label{fig:nlp-analysis-syntax}
\end{figure*}

\subsection{成分结构与解析}
\label{sec:nlp-analysis-constituency}

成分可以写为$(i,j,\ell)$，表示覆盖第$i$个边界到第$j$个边界之间词语的短语，类别为$\ell$。本小节的区间均按词边界编号，不是字符偏移；输出时再经词跨度表还原原文。树中成分必须嵌套或互不相交，不能同时选择$[0,2)$和$[1,3)$这样的交叉区间。

“新／学生／代表”有两种可解释的分组：“新学生”的代表，或新任的“学生代表”。它们的词序相同，分别对应$[[\text{新 学生}]\ \text{代表}]$与$[\text{新}\ [\text{学生 代表}]]$。一个模型可以给两个内部跨度都很高的分数，但合法树不能同时保留它们。本节固定这三个词，分别用文法、神经跨度和动作构造候选，说明评分如何与结构约束配合。

\subsubsection{PCFG与CKY图表解析}
\label{sec:nlp-analysis-pcfg-cky}

\term{概率上下文无关文法}（probabilistic context-free grammar, PCFG）给每条文法规则一个条件概率：同一个左部非终结符的所有展开概率之和为1。一棵树的概率为所用规则概率的乘积，这相当于假定一次展开只依赖该节点的类别，不直接依赖整句其他位置。条件分解的概率观点可回看前卷概率图模型；这里另外建立文法规则和解析算法，不能把一般有向模型直接当作现成的句法解析器。

采用只有$A\to BC$与$A\to w$两类规则、无空产生式和非终结符单分支规则的形式。参考树可以提供规则计数，由
$p(A\to\beta)=N(A\to\beta)/\sum_{\gamma}N(A\to\gamma)$
估计参数；需要平滑时保持同一左部归一化。本例把根类别固定为NP，文法为
\[
\begin{array}{lll}
 \mathrm{NP}\to\mathrm A\ \mathrm{NP}&0.4&
 \mathrm{NP}\to\mathrm{NP}\ \mathrm N\quad0.3\\
 \mathrm{NP}\to\mathrm N\ \mathrm N&0.2&
 \mathrm{NP}\to\text{学生}\quad0.1\\
 \mathrm A\to\text{新}&1&
 \mathrm N\to\text{学生}\quad0.6\\
 &&\mathrm N\to\text{代表}\quad0.4
\end{array}
\]
这是一套只用于手算的微型文法，不宣称覆盖中文。NP规则概率和为1，N规则概率和也为1；无规则的组合不是低概率树，而是不在候选集合中。

\term{CKY图表解析}把每个“区间＋根类别”作为可复用状态。令$V[A,i,j]$为类别$A$覆盖$[i,j)$时的最优子树概率。长度1的初始化为
\[
 V[A,i,i+1]=p(A\to w_{i+1}),
\]
没有相应词汇规则时为0。长度至少2时，根若使用$A\to BC$，左右子树必在某个边界$k$分开，于是
\begin{equation}
 V[A,i,j]=
 \max_{\substack{A\to BC\\i<k<j}}
 p(A\to BC)V[B,i,k]V[C,k,j].
 \label{eq:nlp-analysis-pcfg-cky}
\end{equation}
同时保存取得最大值的$(B,C,k)$及规则身份。按区间长度从小到大填表，每个组合所需的两个子区间都已算好。任意合法二叉树都存在这样的顶层分割；若其中一棵子树不是同一根类别、同一区间下的最优树，替换后会提高整树概率，这就解释了最优子结构为何成立。

表~\ref{tab:nlp-analysis-cky}列出所有非零条目。长度2的“学生代表”有两条推导，必须先比较它们，不能只保留第一次找到的规则。

\begin{table}[htbp]
  \centering
  \begin{tabularx}{\linewidth}{@{}l X l@{}}
    \toprule
    区间&非零类别及计算&最优回溯\\
    \midrule
    $[0,1)$&$\mathrm A:1$&新\\
    $[1,2)$&$\mathrm N:0.6,\ \mathrm{NP}:0.1$&学生\\
    $[2,3)$&$\mathrm N:0.4$&代表\\
    $[0,2)$&$\mathrm{NP}:0.4\cdot1\cdot0.1=0.04$&$\mathrm A,\mathrm{NP},k=1$\\
    $[1,3)$&$\mathrm{NP}:\max(0.2\cdot0.6\cdot0.4,\,
      0.3\cdot0.1\cdot0.4)=0.048$&$\mathrm N,\mathrm N,k=2$\\
    $[0,3)$&$\mathrm{NP}:\max(0.4\cdot1\cdot0.048,\,
      0.3\cdot0.04\cdot0.4)=0.0192$&$\mathrm A,\mathrm{NP},k=1$\\
    \bottomrule
  \end{tabularx}
  \caption{“新／学生／代表”的完整非零CKY图表。未列出的条目为0。根的另一种顶层分组概率为0.0048；右分支内部另有0.012的候选推导，但其最优值为0.048。}
  \label{tab:nlp-analysis-cky}
\end{table}

从根$V[\mathrm{NP},0,3]$读取$(\mathrm A,\mathrm{NP},1)$，再读取右子树的$(\mathrm N,\mathrm N,2)$，得到
\[
 [\mathrm{NP}\ [\mathrm A\ \text{新}]
       [\mathrm{NP}\ [\mathrm N\ \text{学生}]
                    [\mathrm N\ \text{代表}]]].
\]
其概率为0.0192，左分组为$0.3\cdot0.04\cdot0.4=0.0048$。原文“新学生代表”中三个词的字符跨度依次为$[0,1)$、$[1,3)$、$[3,5)$，因此右侧NP恢复为“学生代表”$[1,5)$。概率优势只反映本例文法的偏好；增加前文“刚入学”可能支持另一解释，而这个PCFG并没有直接读取那条背景。

实际文法若有$A\to BCD$，可引入专用辅助类别$A_\star$，改成$A\to B A_\star$、$A_\star\to CD$，把原概率放在第一条规则、第二条设为1，回溯后删除辅助节点。辅助类别必须记录原规则身份，否则可能引入原文法不存在的组合。单分支链$A\to B\to C$可先做带回溯信息的单分支闭包，或把链折叠成复合类别；恢复时展开链。空产生式也需单独消除或改写递推，本式不覆盖它们。

若根值为0，返回该文法下解析失败。对长句使用对数概率，将乘积变成和、0变成$-\infty$，避免下溢。二叉规则集合为$R_2$、非终结符集合为$N$时，直接实现时间为$O(n^3|R_2|)$，图表和回溯空间为$O(n^2|N|)$。把最大值换成求和得到的是内部概率，不再是最优树；要输出树仍需独立的最优指针或采样过程。

\subsubsection{神经跨度成分解析}
\label{sec:nlp-analysis-neural-cky}

PCFG局部规则概率难以直接利用远处词语。神经跨度解析先编码整句，再为每个跨度和类别评分，但仍以合法树为输出。Kitaev与Klein的解析器保留图表解码，并用自注意力编码产生跨度得分；它没有用编码器替代树搜索\scite{nlp-intro-kitaev2018}

将词位置编码拆为用于左右端点的两个子向量，形成边界表示$\bm f_i$。一种与该路线一致的写法是把左右边界差分拼接为$\bm r_{i,j}$，再计算
\[
 \bm s_{i,j}=\bm W_2
  \operatorname{ReLU}\!\left(
   \operatorname{LayerNorm}(\bm W_1\bm r_{i,j}+\bm b_1)
  \right)+\bm b_2.
\]
其中$\bm s_{i,j}$的每个分量对应一个类别。端点索引必须配合句首、句末符号，不应把词的右端表示误当成它的左端。原方法用拆分后的上下文表示构造这类差分，通用注意力计算见第~\ref{chap:transformer}章。

令树得分为$S(T)=\sum_{(i,j,\ell)\in T}s(i,j,\ell)$。单分支链折叠成一个复合标签，多分支树的二叉化辅助节点使用空标签$\varnothing$，并固定$s(i,j,\varnothing)=0$。这样同一原始树不同二叉化不会仅因辅助节点而得到不同总分。这里的“空标签”表示恢复时删除辅助节点，不表示空字符串或遗漏词。

若类别选择不依赖父标签或兄弟标签，令$q(i,j)=\max_{\ell}s(i,j,\ell)$，则
\begin{align}
 C(i,i+1)&=q(i,i+1),\nonumber\\
 C(i,j)&=q(i,j)+\max_{i<k<j}[C(i,k)+C(k,j)].
 \label{eq:nlp-analysis-neural-cky}
\end{align}
每项保存最佳标签和分割点，根禁止空标签，并限定为任务允许的句子或短语根类别。若评分包含父子类别耦合，就须像式~\eqref{eq:nlp-analysis-pcfg-cky}一样把类别留在状态中；若加入跨任意跨度的相互作用，原递推未必仍能精确解码。

对同一“新／学生／代表”，设所有单词叶分数为0，两个长度2跨度的最佳NP分数分别为5和6，根NP分数为1，其他非空类别低于这些分数。独立选跨度会同时取$[0,2)$与$[1,3)$，得到$5+6+1=12$，但二者交叉。图表计算$C(0,2)=5$、$C(1,3)=6$，根比较$1+0+6=7$和$1+5+0=6$，选择右分组。回溯得到NP“学生代表”而不保留NP“新学生”，再删除空标签辅助节点并展开复合单分支标签，才是提交评价的树。

标注树不仅给局部类别，也规定哪些跨度能同时成立。该代表方法采用结构间隔损失
\begin{equation}
 L=\max_T[S(T)+\Delta(T,T^\star)]-S(T^\star),
 \label{eq:nlp-analysis-parse-hinge}
\end{equation}
其中$\Delta$按带标签跨度的差异计代价，且候选包含$T^\star$，所以损失非负。损失可按跨度分解时，可在CKY局部分数中加入代价，找到最违反间隔的树，再提高参考跨度相对其竞争跨度的分数。不能把这个训练目标写成原论文未使用的逐词分类目标。

评价将输出树转成带标签跨度集合，计算精确率、召回率和F1，并固定是否排除根、词性节点、标点及辅助节点。本例若仅比较两个长度2候选且排除根，预测右NP而参考左NP，二者没有匹配成分。树合法只保证结构自洽；长句、深嵌套、未见类别链和领域词仍可能使分数指向错误树。

\subsubsection{Earley文法解析}
\label{sec:nlp-analysis-earley}

二叉化便于CKY组合，但文法作者未必愿意把每条规则改成二叉形式。\term{Earley解析}用带点项目记录“一条规则已经读到哪里”，可直接处理一般上下文无关规则。位于图表位置$j$的项目$[A\to\alpha\mathbin{\cdot}\beta,i]$表示从$i$开始，已经匹配了$\alpha$，仍等待$\beta$；点不是一个输入词元。

以$[S'\to\mathbin{\cdot}\mathrm{NP},0]$初始化位置0。若点后是非终结符$B$，预测操作加入$[B\to\mathbin{\cdot}\gamma,j]$；若点后终结符恰为下一词，扫描操作将点前移一格并放入位置$j+1$；若项目完成为$[B\to\gamma\mathbin{\cdot},i]$，完成操作回到位置$i$，寻找等待$B$的项目，并把它们的点向前移至当前位置$j$。每个位置反复处理新增项目直至闭包，重复项目按规则、点位置、起点和当前位置去重。

沿用上一文法和三个词，位置0预测$\mathrm{NP}\to\mathbin{\cdot}\mathrm A\ \mathrm{NP}$及$\mathrm A\to\mathbin{\cdot}\text{新}$。扫描“新”后，位置1有
$[\mathrm A\to\text{新}\mathbin{\cdot},0]$和
$[\mathrm{NP}\to\mathrm A\mathbin{\cdot}\mathrm{NP},0]$，于是预测从1开始的NP。扫描“学生”后，位置2出现$[\mathrm N\to\text{学生}\mathbin{\cdot},1]$、
$[\mathrm{NP}\to\text{学生}\mathbin{\cdot},1]$；
前者推进$[\mathrm{NP}\to\mathrm N\mathbin{\cdot}\mathrm N,1]$，后者既能完成“新学生”，又能等待后面的N。

扫描“代表”后，位置3完成$[\mathrm{NP}\to\mathrm N\ \mathrm N\mathbin{\cdot},1]$，继而完成$[\mathrm{NP}\to\mathrm A\ \mathrm{NP}\mathbin{\cdot},0]$，最终出现$[S'\to\mathrm{NP}\mathbin{\cdot},0]$。左分组也能完成，不能看到第一个接受项目就宣布它是目标树。为各次扫描和完成保存来源项目及子树指针，形成共享解析森林；按原PCFG的规则乘积评分，就仍选择0.0192的右分组，而不是0.0048的左分组。

本例无空规则与单分支循环，森林上的概率比较可直接自底向上完成。一般实现需要处理可空类别、左递归的项目去重及有循环推导的森林，不能简单递归展开全部路径。Earley免除了文法二叉化要求，但不会免除评分、歧义处理和树恢复；只有接受状态的识别器还没有交付解析树。

\subsubsection{移进--归约成分解析}
\label{sec:nlp-analysis-shift-reduce}

图表同时保留许多区间，也可以按动作逐步构树。设状态包含栈、未读词缓冲区及已建树片段。本文采用一个明确的教学动作系统：$\operatorname{SHIFT}$把下一个带词汇类别的叶入栈；$\operatorname{REDUCE}(A)$把栈顶两个相邻片段按允许规则合成根为$A$的片段；$\operatorname{UNARY}(A)$按允许的单分支规则改包一层。二叉化标记和单分支许可集合预先固定，禁止无限重复套用单分支。

对右分组，“新／学生／代表”的参考动作是
\[
\begin{gathered}
 \operatorname{SHIFT}(\mathrm A,\text{新}),\
 \operatorname{SHIFT}(\mathrm N,\text{学生}),\
 \operatorname{SHIFT}(\mathrm N,\text{代表}),\\
 \operatorname{REDUCE}(\mathrm{NP}),\
 \operatorname{REDUCE}(\mathrm{NP}).
\end{gathered}
\]
前三步栈中有三个叶；第四步按$\mathrm{NP}\to\mathrm N\ \mathrm N$合并后两个，此时栈为
\[
 [\mathrm A(\text{新}),\mathrm{NP}(\text{学生代表})].
\]
第五步按$\mathrm{NP}\to\mathrm A\ \mathrm{NP}$组成根。缓冲区空且栈只剩一个允许的根，才满足完成条件。

训练可从参考树的后序遍历生成合法动作序列，用当前栈顶、缓冲区和句子表示预测动作，最小化参考动作的负对数概率。若保留多个等价二叉化，就需规定训练选哪条轨迹或怎样合并其监督。设读完“新／学生”后，“继续移进”和“立即组合”的局部概率为$.45,.55$，而两条后续完成路径条件概率分别为$.9,.4$，完整概率为$.405$与$.22$。贪心选错早期动作后不能凭后续局部高分拆开已经组合的树；束宽2有机会保留两条并选回右分组，有限束宽仍不保证全局最优。

这个动作系统只能组合相邻栈片段，因而产生连续成分树。恢复时删除二叉化辅助节点、展开复合类别，评价口径与图表法一致。动作正确执行与语境解释正确也须分开：一条完整合法轨迹仍可能选了错误的歧义分组。

\subsection{依存结构与解析}
\label{sec:nlp-analysis-dependency}

基本依存分析以词$1,\ldots,n$和人工根0为节点，边$h\to d$表示中心词$h$支配依存词$d$，边上另有关系标签。本文要求每个真实词恰有一个中心词，人工根没有入边，所有真实词可从根到达，且人工根恰好连接一个真实根词。\term{投射性}（projectivity）进一步要求词及其后代占据连续区间；把人工根放在词序最左侧时，可用依存边无交叉来检查相应树结构。

UD v2提供跨语言共享的词法和依存标注框架，但分词、自由词序和语言特性仍影响候选与评价；完整来源见本章末的文献说明。
无标签依存准确率（unlabeled attachment score, UAS）检查中心词是否正确，带标签依存准确率（labeled attachment score, LAS）还要求关系标签正确。本章手算使用给定分词、计入指向真实根词的根弧，不将人工根作为待预测词；示例没有标点词。实际评测须另行声明是否排除标点，并按数据集指定的标点识别规则处理，不能在系统间改变分母。增强依存则允许额外边和空节点，不能用同一棵基本树的检查条件验收。

\subsubsection{arc-standard与arc-eager转移解析}
\label{sec:nlp-analysis-arc-transition}

两种动作系统都维护栈$\sigma$、缓冲区$\beta$和边集合$E$，从$([0],[1,\ldots,n],\varnothing)$开始，但连边时机不同。arc-standard在栈顶两个词之间连边：若栈末为$i,j$，左弧建立$j\to i$并移除$i$，右弧建立$i\to j$并移除$j$；移进则把缓冲区首词压栈。被移除的词必须尚无中心词，不能移除人工根。参考轨迹还须等被移除节点的参考子节点全部处理完，避免丢失未来边。

对“小王$_1$／阅读$_2$／通知$_3$”，完整轨迹为：移进1，移进2，左弧$nsubj$建立$2\to1$，移进3，右弧$obj$建立$2\to3$，右弧$root$建立$0\to2$。栈依次为
\[
 [0]\to[0,1]\to[0,1,2]\to[0,2]
 \to[0,2,3]\to[0,2]\to[0].
\]
本例根弧只在缓冲区空、唯一子树覆盖全句时建立，防止人工根得到多个子节点。

arc-eager在栈顶$s$与缓冲区首词$b$之间连边：左弧建立$b\to s$并弹出尚无中心词的$s$；右弧建立$s\to b$并把$b$移入栈；归约只弹出已有中心词的栈顶；移进不连边。对同一句，轨迹为移进1，左弧$2\to1$，右弧$0\to2$，右弧$2\to3$，归约3，归约2。它在读入“阅读”时已经连上人工根，而arc-standard等其子树完成后才连接。需要额外禁止第二条根弧，且缓冲区空时若栈中还残留无中心词节点，不能把它们无条件归约后宣称完成。

参考依存树通过确定的oracle转换成动作监督，当前状态的特征或神经表示产生合法动作logit，训练用交叉熵。例如三个合法候选logit为$2,1,0$，概率约为$.665,.245,.090$，贪心选第一项；束搜索累加动作对数概率保留多个状态。屏蔽动作只保证特定局部前提，不证明尚能到达参考树；一旦错误弹出节点，涉及它的后续边可能不可达。没有交换等扩展动作时，这两套系统针对投射结构，不能声称已经覆盖任意非投射树。

\subsubsection{Eisner投射树解码}
\label{sec:nlp-analysis-eisner}

动作模型逐步选择构树路径，图方法则先给候选边评分，再求整棵树的总分$\sum_{(h,d)\in E}s(h,d)$。逐词选最高入边可能产生环；即使打断环，还可能违反投射性。Eisner算法用连续区间上的四类状态解决这个约束问题，其图方法背景见原始依存解析工作和最大生成树表述，完整来源列于本章末的文献说明。

本小小节用闭区间$[i,j]$表示从词$i$到词$j$，与成分解析的边界区间区别开。$C^R_{i,j}$表示以左端词$i$为根的完整区间，$C^L_{i,j}$表示以右端词$j$为根的完整区间；$I^R_{i,j}$和$I^L_{i,j}$分别显式包含边$i\to j$和$j\to i$，但端点一侧的子树还需与另一个完整区间组合，因此称为不完整状态。初始化
$C^R_{i,i}=C^L_{i,i}=0$，其余尚未构造项为$-\infty$。

先把两个不重叠的完整区间连成一条端点边：
\begin{align}
 I^R_{i,j}&=s(i,j)+
  \max_{i\leq k<j}(C^R_{i,k}+C^L_{k+1,j}),\nonumber\\
 I^L_{i,j}&=s(j,i)+
  \max_{i\leq k<j}(C^R_{i,k}+C^L_{k+1,j}).
 \label{eq:nlp-analysis-eisner-incomplete}
\end{align}
再补齐共享端点外侧的子结构：
\begin{align}
 C^R_{i,j}&=\max_{i<k\leq j}
            (I^R_{i,k}+C^R_{k,j}),\nonumber\\
 C^L_{i,j}&=\max_{i\leq k<j}
            (C^L_{i,k}+I^L_{k,j}).
 \label{eq:nlp-analysis-eisner-complete}
\end{align}
两项共享的端点不是重复添加的边，组合只累加各自已建立的边分数。按区间长度先算不完整项、后算完整项，记录每次最大值的分割点和状态类型。为了严格保证一条根弧，在真实词的图表完成后计算
\[
 \max_{1\leq r\leq n}
 [s(0,r)+C^L_{1,r}+C^R_{r,n}],
\]
保存最优根$r$。这相当于把一棵投射树分为根词左侧、右侧两部分，而不是允许人工根自由连接多个子树。

给“小王／阅读／通知”构造如下全部有限边分数：
\[
 s(0,2)=5,\quad s(2,1)=4,\quad s(2,3)=3,\quad
 s(3,1)=5,\quad s(1,3)=6.
\]
其余边禁止，包括自环和指向人工根的边。局部最高入边为$3\to1$、$0\to2$、$1\to3$，总分16，却在1和3之间成环。长度2图表得到$I^L_{1,2}=C^L_{1,2}=4$、
$I^R_{2,3}=C^R_{2,3}=3$，相反方向都为$-\infty$。虽然边$1\to3$有6分，式~\eqref{eq:nlp-analysis-eisner-incomplete}中两个分割都需要禁止的$1\to2$或$3\to2$，所以相应不完整项不可达。

唯一允许的根为2，完整得分为$5+4+3=12$。回溯左半区间取$2\to1$，右半区间取$2\to3$，加根弧$0\to2$，正好恢复图~\ref{fig:nlp-analysis-syntax}的无标签树。标签可在边选定后另行最大化。时间为$O(n^3)$，图表与指针为$O(n^2)$；这个精确性只针对可加的边分数和投射树空间，并不保证语言分析正确。

\subsubsection{Chu--Liu/Edmonds非投射树解码}
\label{sec:nlp-analysis-mst}

自由词序和远距离依赖可能需要非投射树。最大权有向生成树仍要求单中心词、无环和根可达，只是取消投射性。\term{Chu--Liu/Edmonds算法}先选每个非根节点的最大入边；若无环，便已达到逐节点上界。若有环，将环收缩为一个节点，重新计算进出环的边分数，递归求解后展开。

设选出的环为$Q$，环内每个节点$v$的已选前驱为$a(v)$，环内总分为$w(Q)=\sum_{v\in Q}s(a(v),v)$。一条外部边$u\to v$进入环，会替换掉原入边$a(v)\to v$，其收缩得分可写为
\begin{equation}
 s'(u,Q)=\max_{v\in Q}
   [s(u,v)-s(a(v),v)+w(Q)].
 \label{eq:nlp-analysis-mst-contract}
\end{equation}
出环边取$\max_{v\in Q}s(v,u)$。两种最大值都要保存原始端点；展开时保留被选择的外部入边，删去它所替换的那一条环边，补回其他环边。只做收缩而不记录端点，无法恢复原词之间的依存。

沿用上一组得分，环$Q=\{1,3\}$的总分为$5+6=11$。从2进入1的收缩分数为$4-5+11=10$，进入3为$3-6+11=8$。收缩图选择$0\to2$的5分和$2\to Q$的10分；展开时用$2\to1$替换$3\to1$，保留$1\to3$，得到总分15的树。该树无环，却因$0\to2$与$1\to3$交叉而非投射；Eisner排除它，故最高分只有12。这是允许结构改变导致的结果差异，不能据此说15分的语言解释更好。

一般生成树算法可能让人工根连接多个真实词。若规范只允许一个根词，可枚举$r$，仅保留$0\to r$这一条人工根出边，并禁止其他词指向$r$，分别求树后选最高分；这会增加计算量，但给出清楚的单根保证。本例已经只允许$0\to2$。若任何词没有可用入边，或不存在从根覆盖全部词的结构，返回失败。朴素递归收缩可达$O(n^3)$，专门的数据结构能降低稠密图复杂度；不能把一种优化实现的复杂度无条件归给所有收缩代码。

\subsubsection{双仿射依存解析}
\label{sec:nlp-analysis-biaffine}

上述解码器假定边分数已经给定。双仿射解析把“谁适合作中心词”和“谁适合作依存词”投影到不同表示，再为二者配对评分。设$\bm h_i^{\mathrm{head}},\bm h_j^{\mathrm{dep}}\in\mathbb R^d$由上下文编码和两个前馈投影得到，则
\begin{equation}
 s(i,j)=
 (\bm h_i^{\mathrm{head}})^\top\bm U\bm h_j^{\mathrm{dep}}
 +\bm a^\top\bm h_i^{\mathrm{head}}
 +\bm b^\top\bm h_j^{\mathrm{dep}}+c.
 \label{eq:nlp-analysis-biaffine}
\end{equation}
双线性项比较两端的相容性，仿射项允许中心词或依存词各有偏好。Dozat与Manning的代表方法分别构造弧与关系的表示及评分，不能把一个关系logit当成已经选定的中心词\scite{nlp-analysis-dozat2017}

对每个依存词$j$，在允许中心词上做softmax，以标注中心词$i^\star$训练$-\log p(i^\star\mid j,x)$；关系头在参考边两端上训练正确关系的交叉熵。推断时先将边分数送入投射或非投射树解码，再在选中的边上恢复关系。若标签分数也进入树目标，应明确先计算$\max_\ell s(i,j,\ell)$并保存标签，再解码相应树；这与只按弧分数建树是两种设置。

取$d=2$、$\bm U=\bm I$且仿射项为0，若“阅读”的中心词表示为$(1,0)$、“小王”的依存表示为$(4,0)$、“通知”的依存表示为$(3,0)$，则两条边分数为4和3。这正好可以接入前面的投射例子。选定$2\to1$后，关系头若给nsubj、obj、其他的logit为$2,0,-1$，选择nsubj；对$2\to3$则可选obj。恢复带标签边及词的原文跨度后，才交付完整解析。

若两个中心词都正确但把$2\to1$标成obj，含根的三条边UAS为1，LAS为$2/3$；如果中心词本身错了，正确预测一个关系名也不能弥补。上下文编码器、边评分、树解码和标签分类的错误应分开记录，解析指标提高是否改善抽取或生成，还要在下游任务中检验。

\subsubsection{增强依存图预测}
\label{sec:nlp-analysis-enhanced}

基本树常常只保留一个句法中心词，但共享论元和省略关系可能需要额外连接。\term{增强依存}（enhanced dependencies）按所选规范在词节点以及允许的空节点上增加关系；UD的增强表示允许多入边、空节点，某些结构还可能出现环。它不是把基本树限制放宽后任意加边，哪些边有意义仍由标注规范规定。

例如补充教学句“小王购买并阅读通知”，基本树可有“购买”到“小王”的nsubj、“购买”到“阅读”的conj、“阅读”到“通知”的obj。增强表示可传播共享主语，增加“阅读”到“小王”的nsubj；在语义和规范允许时，还可传播共享宾语，增加“购买”到“通知”的obj。“小王”因此有两个入边，却不是违反增强表示的错误。

一种实现先保留基本树，再对允许的节点对和关系$\ell$计算
$p_{i,j,\ell}=\operatorname{sigmoid}(g_\ell(i,j,x))$，以增强边多热标签训练二元交叉熵。假设新增两条边的概率为$.85,.80$，反向错误候选为$.15$，验证集确定阈值$.5$时，保留前两条并记录标签。边阈值是本例的恢复方案，不是UD强制采用的预测算法。多标签是否允许同一节点对出现多个关系，也需依据实际规范检查。

对“甲买书，乙报纸”这类省略结构，还需要先预测第二分句的缺失谓词位置及空节点身份，再把“乙”“报纸”的边接到该节点；不能只把一个未出现在原文的“买”加入字符跨度列表。应分别检查空节点锚点、边端点是否存在、关系是否合法及增强图所要求的连通性，不再施加基本树的单入边或统一无环条件。增强边和标签F1与UAS、LAS的分母、对象不同，应分别报告。

\section{语义分析}
\label{sec:nlp-analysis-semantics}

句法可以判定“小王”是“交给”的主语，却不会自动给出“谁接收了什么”。\term{语义分析}还要约定意义以什么形式交付：一个词义编号、一组谓词论元、一张概念图，或一个能在指定环境中求值的表达式。它们回答的问题不同，不存在只要预测一种就完成其他所有分析的保证。表~\ref{tab:nlp-analysis-semantic-outputs}汇总了本节的五类对象；后面的例子逐步说明这些输出怎样形成。

\begin{table*}[htbp]
  \centering
  \small
  \setlength{\tabcolsep}{3pt}
  \begin{tabular}{@{}p{15mm}p{24mm}p{32mm}p{24mm}p{27mm}p{28mm}@{}}
    \toprule
    预测对象&教学输入&具体输出&结构检查&意义检查&未覆盖信息\\
    \midrule
    词义&river bank&bank的河岸义编号&编号属于候选清单&所选释义符合当前用法&未收录的新义\\
    语义角色&小王把书交给小李&交给：施事小王、受事书、接受者小李&跨度与角色符合规范&论元与当前谓词对应&事件是否真实发生\\
    AMR&少年想阅读&want-01、read-01及共享施事节点&引用存在、关系与图合法&意愿与阅读的嵌套正确&实际是否阅读\\
    逻辑式&并非所有来宾已报名&$\neg\forall x(G(x)\Rightarrow R(x))$&变量作用域与类型正确&区别于所有人均未报名&事实库外的信息\\
    SQL&已报名校外来宾有几人&连接来宾与报名表后计数&列、连接、聚合合法&条件与计数对象准确&数据库未记录的状态\\
    \bottomrule
  \end{tabular}
  \caption{语义分析的产出与核验。五行的结构对象不同；向量可以参与评分，但不是这些任务的最终交付结果。逻辑式和SQL的求值都依赖明确给定的环境。}
  \label{tab:nlp-analysis-semantic-outputs}
\end{table*}

\subsection{词义消歧与词汇意义}
\label{sec:nlp-analysis-wsd}

“bank”在“deposit money at the bank”和“sit by the river bank”中具有不同用法。\term{词义消歧}（word sense disambiguation, WSD）先按基本形式检索候选，再从固定词义清单中选择身份；这里设清单只有金融机构义$F$和河岸义$R$。这些是教学编号，不冒充WordNet的正式sense key。基本形式、词性限制和清单版本都影响候选范围，不能把不同清单下的准确率直接比较。

固定清单时，最简单的监督基线是用上下文特征预测词义类别；训练标签来自逐词义标注。若某词训练样本全是金融义，恒选$F$也可能取得很高总准确率，却不能处理少见用法。释义或语义网络能提供额外依据，但仍需为未收录用法保留“未知／待补录”。词义正确与翻译、检索正确应分别检查；选择$F$也不等于已经输出bank与其他词之间全部同义、上下位或反义关系。

\subsubsection{Lesk与简化Lesk释义重叠}
\label{sec:nlp-analysis-lesk}

词典释义也包含可以与语境比较的词。Lesk路线在各词候选释义之间寻找重叠；简化Lesk只比较目标词的候选释义与原文语境，因而不必同时为每个邻词决定词义。设规范化后语境词集合为$C$、候选释义词集合为$G_s$，一种简单分数是
\[
  o(s)=|C\cap G_s|,\qquad \hat s=\argmax_{s\in\{F,R\}}o(s).
\]
规范化过程应固定小写化、词形还原、停用词表及目标词是否删除。本例删去目标词与停用词，将两条释义设为
\[
\begin{aligned}
 G_F&=\{\code{money},\code{deposit},\code{institution}\},\\
 G_R&=\{\code{land},\code{river},\code{edge}\}.
\end{aligned}
\]
两句语境分别得到$C_1=\{\code{deposit},\code{money}\}$和$C_2=\{\code{sit},\code{river}\}$，分数依次为$(2,0)$与$(0,1)$，恢复$F$与$R$。

若采用原始的释义间比较，可为上下文词也枚举释义，最大化$\sum_{i<j}|G_{s_i}\cap G_{s_j}|$。例如“river”只有释义$\{\code{water,land,flow}\}$，与bank两义的重叠为0与1，仍选$R$；若上下文词都多义，候选组合数会随词数指数增长，需要限制窗口、候选或采用近似搜索。这里给出的是集合重叠实现，计重复词、短语重叠及长度归一化都是不同变体，必须说明再比较。

较长释义有更多偶然命中机会，可改用Jaccard分数$|C\cap G_s|/|C\cup G_s|$。第一句此时给$F$打$2/3$分而非2分。语境“withdraw cash”与上述两条释义都无字面重叠，并不意味着两个词义都不成立；它暴露的是词汇覆盖不足。零分平局应保留歧义或按事先约定的词频回退，不能把平局打破方式当作语义证据。

\subsubsection{GlossBERT语境--释义判定}
\label{sec:nlp-analysis-glossbert}

同义表达缺少共同词时，可以让模型联合读取语境和释义，判断这个候选是否适用。GlossBERT按目标基本形式取出全部候选，将每个“语境--释义”对变成二分类样本；正确词义的标签为yes，其余为no。其目标位置标记版本还强调句中是哪一次词语出现，避免同句多个相同词面混淆\scite{huang-etal-2019-glossbert}

输入可写成\code{[CLS] context [SEP] gloss [SEP]}，并在目标词两侧加约定标记。句对编码得到$\bm h_s$，二分类头输出$p_s=\operatorname{sigmoid}(\bm w^\top\bm h_s+b)$，训练累加正负候选的二元交叉熵。假设河岸句对两个候选的logit为$-1,2$，则yes概率约为$.269,.881$；推断取最大值对应的编号$R$。这两个概率来自分别进行的二分类，不要求相加为1。输出必须保存候选编号及清单版本，不能只返回释义字符串让后续系统猜身份。

如果句中是“the bank financed repairs to the river bank”，两个目标的位置会让同一上下文产生不同预测。释义极接近时也需保留目标位置和周围句法信息；联合编码允许两段文字逐词交互，但并不保证学到正确区分。检索步骤若只返回金融义，评分再高也不能证明答案正确；候选遗漏、正确候选低分和清单没有所需新义应分别记录。BERT的编码机制见第~\ref{chap:transformer}章，这里的新增监督是词义与候选的配对关系。

\subsubsection{释义双编码检索}
\label{sec:nlp-analysis-gloss-retrieval}

联合编码需要对每条释义重读语境。候选很多时，可以分别编码语境中的目标词和释义，再用向量内积比较。Blevins与Zettlemoyer的双编码方法联合训练两个编码器：目标词表示来自上下文位置，释义表示来自其序列汇聚位置，训练对当前词的候选词义做softmax交叉熵\scite{blevins-zettlemoyer-2020-moving}

设$\bm q=f(x,\text{目标跨度})$，$\bm g_s=g(\text{释义}_s)$，则
\[
 p(s\mid x)=\frac{\exp(\bm q^\top\bm g_s)}
                  {\sum_{t\in S(w)}\exp(\bm q^\top\bm g_t)}.
\]
多个模型词元对应一个目标词时，须先按位置映射汇聚。用$\bm g_F=(1,0)$、$\bm g_R=(0,1)$以及河岸语境$\bm q=(0,2)$，两候选得分为0和2，概率约为$.119,.881$；检索返回第二个向量的身份$R$，不是向量本身。金融语境$\bm q=(2,0)$则给相反排序。

同一候选清单下，释义向量可预计算，语境只编码一次，再付出$O(|S(w)|d)$内积成本；大清单可用近似检索，但需另查候选召回。它放弃了每次配对时的细粒度交叉注意力，计算代价与交互方式都不同于GlossBERT。最高内积无论多大都不能保证清单覆盖未知用法；拒绝阈值或分数间隔需用含未知样本的验证集确定，不能直接沿用本例的任意数值。

\subsubsection{词汇语义图上的个性化PageRank}
\label{sec:nlp-analysis-wsd-ppr}

释义之外，词义之间的语义关系也能传播语境信息。把词义作为节点、词汇语义关系作为边，并把语境词映射到其候选节点，便可从语境附近开始随机游走。个性化PageRank用非均匀重启分布表达这种偏置；在词义消歧中，使用所有节点等概率重启的全图PageRank不能代替语境条件排序\scite{agirre-soroa-2009-personalizing}

令$\bm P$为按列归一化的转移矩阵，$\bm v$为非负且和为1的语境分布，迭代
\begin{equation}
 \bm r^{(t+1)}=\rho\bm P\bm r^{(t)}+(1-\rho)\bm v,
 \qquad 0<\rho<1.
 \label{eq:nlp-analysis-personalized-pagerank}
\end{equation}
没有出边的节点按约定转向$\bm v$，使每列仍和为1。可从$\bm r^{(0)}=\bm v$开始，直到相邻两次的$L_1$差小于阈值。对固定随机矩阵，$\rho<1$带来收缩和唯一不动点；稀疏重启分布并不保证每个节点最后都有正概率，不能混同于全支撑重启的不可约性。

构造四节点图，顺序为金融义$F$、贷款义$L$、河岸义$R$、河流义$W$，只有无向边$F$--$L$及$R$--$W$。贷款语境映射成$\bm v=(0,.9,0,.1)$，取$\rho=.5$。一次迭代得$(.45,.45,.05,.05)$，第二次为$(.225,.675,.025,.075)$；极限为$(.3,.6,1/30,1/15)$。只在目标词候选$F,R$中比较，选$F$。河流语境把$.9$和$.1$交换，极限相应交换两组，选$R$；均匀重启则四节点均为$.25$，无法区分两种语境。

真实词汇图中，一个语境词可能对应多个节点，可将它的重启质量分配给全部候选；是否排除目标词自身也须固定。本例的$.9/.1$是给定权重，不是游走自行理解了语境。错误的词形映射、图中缺边和过宽的语境都会影响排序，且关系传播只能选择已有词义，不能创造清单外的新义。

\subsection{谓词、论元与语义角色}
\label{sec:nlp-analysis-srl}

在“小王昨天把一本书交给小李”中，“交给”描述一个关系，其他片段说明参与者及时间。\term{谓词}（predicate）是这个关系或事件的表达，\term{论元}（argument）是与之对应的参与对象，\term{语义角色}（semantic role）说明对象怎样参与。为便于比较，本节教学规范采用施事、受事、接受者和时间，分别记作A0、A1、A2、TMP；这些解释只针对当前谓词，不能把任何资源的ARG0都无条件等同于施事。

给定分词为“小王／昨天／把／一本书／交给／小李”，原文跨度依次为$[0,2)$、$[2,4)$、$[4,5)$、$[5,8)$、$[8,10)$、$[10,12)$。期望输出为以“交给”$[8,10)$为谓词的四条带角色记录。是否已给定谓词位置，决定系统还需不需要识别“交给”；评价完整的“谓词、论元位置、角色”三元组，才能区分这两种条件。此处恢复语言中的参与关系；按业务事件本体组织事实由第~\ref{chap:nlp-information-extraction}章承担。

\subsubsection{成分候选与角色分类}
\label{sec:nlp-analysis-srl-constituents}

如果已有成分树，可将每个非谓词成分当作候选，用其短语类别、中心词、相对位置及到谓词的树路径预测角色或NONE。例如候选“小王”“昨天”“一本书”“小李”分别有局部表示$\bm h_a$，与谓词表示及路径特征拼接后输入softmax分类器；训练标签由参考论元与候选跨度精确匹配形成，未匹配者标为NONE。参考论元不在树中时，不能为了得到训练标签随意扩大到最近成分。

设每个候选的最佳非空角色分数与NONE分数分别为$(3,0)$、$(2,0)$、$(4,0)$、$(3,0)$，最佳角色依次为A0、TMP、A1、A2；二候选比较下概率为$.953,.881,.982,.953$，恢复四条期望记录。另有父成分“把一本书”，其A1得分2低于“一本书”的4；本例规范不允许同一谓词的论元重叠，故两者不能同时保留。可以最大化相对NONE的总分，并对重叠候选加互斥约束；按跨度排序后用区间动态规划即可实现不重叠选择。

若解析器漏掉“一本书”这个候选，即使分类器对所有现有成分都判断正确，A1仍不可能精确恢复。句法候选召回与角色分类准确率因此必须分开。跨越不连续短语的论元、同一核心角色能否重复及省略论元是否显式标注，由角色资源规定；一个连续成分分类器不天然支持所有这些情况，也不能补入原文没有表达的原因。

\subsubsection{BIO--CRF语义角色标注}
\label{sec:nlp-analysis-srl-crf}

不依赖成分候选时，可以对每个谓词单独给词序列打BIO角色标签。输入加入谓词位置指示，标签为O及各角色的B、I；训练与合法转移使用第~\ref{sec:nlp-analysis-bmes-crf}节的CRF，但掩码改为BIO规范：I-A1只能跟在B-A1或I-A1后，其他角色同理。谓词自身在本例记O，谓词位置由另一个字段保存。

对上面的六个词，目标序列为B-A0、B-TMP、O、B-A1、O、B-A2。令这些位置的目标标签得分均为2，其他合法标签得分为0；额外令第四词I-A1得分为3。逐位置贪心得到的I-A1不能接在第三词O后。CRF解码在合法转移为0时取目标序列，总分12；另一条只将第三词改为B-A1、第四词改为I-A1的合法竞争路径总分11。回溯后依据词到字符的映射，恢复四个角色跨度，而不是输出模型词元编号。

“小王答应小李阅读通知”含“答应”和“阅读”两个谓词。对“答应”，小李可为接受承诺者，阅读通知可为承诺内容；对“阅读”，小王为阅读者、通知为受事，小李不是相同角色。必须生成两个谓词条件序列，不能让一个BIO序列同时承担全部关系。若规范要求在当前句中链接被控制的主语，可从“小王”的显式跨度恢复；若主语完全未表达，普通BIO没有可贴标签的位置。

转移掩码只处理局部边界，并不禁止远处出现第二个B-A0。核心角色唯一性需要扩展状态记录已用角色，或在解码后做显式结构选择；允许重复的时间、地点等附加角色不能机械套用同一限制。一个序列也不能表达重叠论元，需要多层标签或换用跨度表示。

\subsubsection{联合跨度式语义角色标注}
\label{sec:nlp-analysis-srl-span}

跨度方法直接枚举论元区间，并与候选谓词配对。He等的联合模型分别为谓词、论元打分，再为“谓词--论元--角色”打分；以空角色分数0作为参照，对每个候选对训练类别交叉熵。为控制原本$O(n^3|L|)$的配对量，按一元分数保留有限谓词、论元，并限制论元最大宽度\scite{nlp-intro-he2018}

对非空角色$r$，可写
\[
 s(p,a,r)=s_p(p)+s_a(a)+s_r(p,a),\qquad
 s(p,a,\mathrm{NONE})=0.
\]
跨度向量由两端表示、内部注意力汇聚及宽度嵌入构成。训练监督同时指定谓词位置、论元跨度和角色；被当作候选而没有标注关系的对提供NONE监督，不能只用正例训练。裁剪前还要检查训练参考候选的覆盖，否则裁剪掉的正例不会参与关系学习。

对“交给”与“一本书”，设三项为$1,1,2$，A1总分4；与“书”配对时为$1,.5,1.5$，总分3；与“把一本书”配对时总分2。独立选标签可能保留三个重叠A1。本例采用不重叠、A1唯一的约束，比较后保留4分的完整跨度$[5,8)$。这是明确增加的结构解码；局部softmax本身没有消除冲突，原始独立预测路线也不能被描述成自动满足全部角色约束。

若谓词“交给”一元得分过低而被裁掉，所有角色一起漏失；若只剪掉“一本书”，可能留下不完整的“书”；若跨度保留但A2胜出，则是角色错误。分别报告谓词召回、跨度候选召回和最终三元组F1，才能知道扩大候选是否有用。连续跨度假设仍不覆盖不连续论元，最大宽度限制则可能漏掉长论元。

\subsubsection{依存式语义角色标注}
\label{sec:nlp-analysis-srl-dependency}

另一些资源用论元的中心词表示参与对象，此时输出是谓词词节点到论元词节点的带角色边。编码每个词后，可用双仿射或拼接前馈网络计算$s(p,j,r)$，在含NONE的角色类别上训练交叉熵。候选谓词可以给定或另行预测；推断对每个候选边选择角色，并按资源检查允许的重复角色及边组合，不要求所有词组成一棵句法树。

将例句细分为“小王／昨天／把／一／本／书／交给／小李”时，完整受事仍是“一本书”$[5,8)$，中心词却是“书”$[7,8)$。假设“交给”到“书”的A1、A2、NONE分数为$3,0,0$，概率约为$.909,.045,.045$，便恢复中心词角色边；其他边同理得到小王、昨天、小李。这个输出没有声称已经恢复受事的完整字符范围。

若下游需要完整短语，可另用句法子树、边界预测头或标注好的头到跨度映射扩展。“书”的句法子树未必恰好覆盖语义论元，所以扩展也要评价。跨度式和中心词式三元组的分母不同，不能直接比较F1。语义边也不能从句法nsubj、obj简单改名得到：被动句的句法主语可以是受事，同一词还可能参与多个谓词。

\subsection{句子语义图}
\label{sec:nlp-analysis-amr}

多个谓词可能谈到同一个对象。“少年想阅读”中，少年既是“想”的主体，也是所期望的“阅读”的主体。\term{抽象意义表示}（abstract meaning representation, AMR）用概念节点、带标签关系和属性组织这些内容；变量标识节点身份，同一变量可以作为多条关系的端点。共享节点的多入边称为\term{重入}（reentrancy），它不是把相同名字写两遍。

采用英语AMR概念标签说明这句中文教学材料，可用PENMAN式记法写成：
\begin{codeblock}{共享论元的AMR}
(w / want-01
   :ARG0 (b / boy)
   :ARG1 (r / read-01
            :ARG0 b))
\end{codeblock}
斜线后是概念，冒号后是关系，末尾的b引用已建立的节点。图中有三个概念节点及三条关系；原文中只有一次“少年”，图里也只有一个b。AMR抽象掉部分表层形式，却不自动断言阅读实际发生；\code{read-01}处于意愿内容中。概念词表、角色定义、属性类型和图合法性依资源版本约定，本节按连通、带根、有向无环且允许重入的图处理。

Smatch先搜索两张图之间的一对一变量对齐，再比较概念、属性和关系三元组；仅比较变量字母会把相同图误判成不同。对上述图，若只省掉$r\xrightarrow{\mathrm{ARG0}}b$，在不计额外TOP三元组的教学口径下，参考有6个三元组、预测有5个，最佳对齐匹配5个，故精确率1、召回率$5/6$、F1为$10/11$。实际评测应固定工具的TOP、属性及规范化口径。对齐搜索可用整数规划或多次重启的局部搜索；后者可能错过最佳对齐，结构匹配本身也不能证明两句话意义等价\scite{cai-knight-2013-smatch}

\subsubsection{CAMR转移式AMR解析}
\label{sec:nlp-analysis-camr}

CAMR从句法依存结构出发，通过动作改变节点、边和标签。它维护当前跨度图、待处理节点队列及当前节点的待处理子边；初始顺序来自依存树的后序遍历，节点和边的语义标签尚未确定。参考AMR与词语对齐后形成跨度图，再由oracle产生动作监督，分类器根据当前图与上下文选择动作\scite{wang-etal-2015-transition}

对“少年$_1$／想$_2$／阅读$_3$”，初始依存边为$2\to1$与$2\to3$，后序队列为$[1,3,2]$，暂不写人工根。节点1没有子边，执行\code{NEXT-NODE(boy)}并前进；节点3也没有子边，执行\code{NEXT-NODE(read-01)}。处理节点2的第一条边$2\to1$时，执行\code{REENTRANCE(3,ARG0)}，增加$3\to1$而保留原边；该动作不推进当前子边。随后\code{NEXT-EDGE(ARG0)}给$2\to1$标注并推进，\code{NEXT-EDGE(ARG1)}处理$2\to3$，最后\code{NEXT-NODE(want-01)}标注节点2并清空队列。将节点1、2、3分别重命名为b、w、r，即恢复上面的图。

重入动作不能创建自环、重复同一边或破坏所选图约束，也要限制重复使用，以免不前进的动作形成无限轨迹。若当前状态的合法候选“增加阅读者”“直接处理当前边”得分为2和0，softmax给前者约$.881$；训练提高参考动作的条件概率，解码累加整条动作轨迹的对数概率，而不是仅根据某一动作是否合理验收。

其他结构差异用不同动作处理。\code{SWAP}反转当前中心词关系，\code{REATTACH}删除旧边后改接父节点，\code{MERGE}合并多词跨度；\code{REPLACE-HEAD}删除当前中心节点并由子节点继承连接，\code{DELETE-NODE}只在相应前提下删除无子边节点。这些动作不能互相替代，例如增加共享论元必须保留原关系，不能误用改接动作。姓名等折叠子图还需按预处理约定展开。错误句法输入可能限制可恢复目标，正确输入也可能被错误动作破坏，二者要分别定位。

\subsubsection{概念识别与关系图构造}
\label{sec:nlp-analysis-amr-node-edge}

另一条路线先预测概念，再决定概念怎样连接。训练时需要文本与概念的对齐或可学习的对齐机制；从词语或跨度提出零个、一个或多个概念候选，并以概念分类、复制或生成损失训练。它不能简单规定“一个词恰好对应一个概念”，否则多词姓名和不直接对应词面的抽象节点都会被漏掉。

设“少年”“想”“阅读”三个跨度各自给目标概念的概率为$.9,.8,.9$，保留b、w、r。关系分类器对每个有序节点对预测关系或NONE；给$w\to b$的ARG0、$w\to r$的ARG1、$r\to b$的ARG0概率为$.9,.8,.7$，反向$r\to w$候选为$.2$。教学恢复规则取阈值$.5$，前三条进入图，按根预测选择w，随后检查可达性与无环性。若出现环，可在候选图上最大化带约束边分数，或按明确规则删边；仅阈值化没有合法图保证。

这里必须允许b有两条入边。若为“阅读者”另建$b'$且标签也为boy，虽然全部概念名看起来正确，却表达了两个对象，遗漏了重入。反之，把两个真实不同的人合并成一个节点会造成错误共享。节点身份、概念标签、关系方向及重入边应分别检查；独立节点预测需要保留稳定内部编号，不能以同名概念自动合并。图到文本可以复用这些编号关系，但怎样把图表达成自然句子由第~\ref{chap:nlp-text-generation}章展开。

\subsubsection{SPRING图线性化与序列到序列解析}
\label{sec:nlp-analysis-spring}

只要序列化过程不丢失邻接与引用信息，也可以把图转成符号序列来训练。SPRING使用预训练编码器--解码器，并比较PENMAN、深度优先和广度优先线性化；其深度优先版本用专用指针符号表示节点变量，避免普通字母既可能是常量又可能是变量的混淆\scite{nlp-analysis-spring2021}

上面的图按深度优先顺序可写成：
\begin{codeblock}{带节点指针的图序列}
( <R0> want-01
  :ARG0 ( <R1> boy )
  :ARG1 ( <R2> read-01 :ARG0 <R1> ) )
\end{codeblock}
首次带概念出现的指针建立节点，后续只出现的\code{<R1>}引用同一节点。这个序列与图可相互恢复；将末项换成\code{(<R3> boy)}则变成另一个少年，不是无害的变量改名。训练用文本--图序列对最小化每步条件交叉熵，目标还包含括号、关系、指针及结束符；AMR符号词表和序列化顺序在训练、解码时须一致。

恢复程序读取左括号并压栈，建立指针到节点的映射，遇关系时记录其源节点，再解析目标子图或指针引用；右括号退栈，结束后解析所有引用并检查唯一概念定义与连通性。本例共建立3个节点、3条边，两个\code{ARG0}的目标都指向节点1。若最后生成未定义的\code{<R9>}，应定位为悬空引用；有些序列化允许前向引用，因此不能一见未定义指针就判错，而应在全序列读完后仍未定义时拒绝。

束搜索比较完整序列概率。假设正确引用\code{<R1>}与新建节点分支的条件概率为$.8,.2$，此前前缀概率相同，则这一步偏向共享节点；后续概率仍会影响全序列排名。SPRING原文还使用轻量后处理处理括号和非法延续，这类修复要记录，不能据可解析性宣称意义已正确。文本到图与图到文本具有对称输入输出接口，但需要相应方向的训练；AMR缺少的时态细节或原始措辞不会因反向生成自动唯一恢复。

\subsection{语义解析与形式表达}
\label{sec:nlp-analysis-formal}

如果目标是计算数量或执行查询，还需要比概念图更明确的类型、作用域和求值规则。\term{语义解析}在这里指把文本映射为指定形式语言中的表达式；结构是否合语法、能否在环境中求值、是否对应原问题，是三层不同判断。下文先用有限事实域说明逻辑式，再用数据库说明SQL，不能把SQL列名的约束套到AMR变量上。

逻辑例子的实体域为$E=\{a,b,c\}$，表示三位不同来宾。给定$G=E$为来宾集合、$R=\{a,b\}$为已报名者，关系$\mathrm{help}$的外延为$\{(a,b),(c,b)\}$。本例采用封闭、完整的有限事实环境：未列出的报名事实为假。现实资料若不完整，缺记录只能得到未知，不能直接移用该封闭世界解释。

数据库例子则额外提供两张教学业务表，不声称原通知已经包含报名状态。\code{people}表的字段为\code{pid}、\code{name}和\code{affiliation}，三行数据分别是
\[
\begin{aligned}
 &(1,\text{小王},\text{external}),\\
 &(2,\text{小李},\text{internal}),\\
 &(3,\text{小赵},\text{external}).
\end{aligned}
\]
\code{registrations}表的字段为\code{event}、\code{pid}和\code{status}，三行为
\[
\begin{split}
 &(\code{ML-0418},1,\code{confirmed}),\\
 &(\code{ML-0418},3,\code{pending}),\\
 &(\code{OTHER},3,\code{confirmed}).
\end{split}
\]
\code{pid}是people的主键，registrations的主键是$(\code{event},\code{pid})$，后者的pid引用前者；OTHER只表示另一个教学活动。状态都非空，同一人在同一活动最多一条记录。问题是“活动ML-0418中已确认报名的校外来宾有几人”，期望值为1。原通知的报名截止条件仍为2025年4月16日18:00前；当前补充表没有提交时间，不能进一步回答是否按时提交。

这类输入假设也是基准之间的重要差别。Spider（\citeyear{yu-etal-2018-spider}）按数据库划分考查跨域模式泛化，Spider 2.0（\citeyear{nlp-analysis-spider2025}）进一步纳入实际工作流中的复杂模式、方言和元信息需求；它们不能仅被理解成更长的同一道SQL题。SParC提供上下文相依的连续查询，CoSQL还涉及数据库对话中的澄清与不可回答情形，四项资料见本章末的文献说明。元信息、历史轮次和可检索示例是否可见，都应进入任务定义；本章讨论语义映射，结果如何组织成答案以及跨步骤执行恢复分别由第~\ref{chap:nlp-question-answering-dialogue}、\ref{chap:nlp-solving-and-execution}章承接。

\subsubsection{类型化文法约束的逻辑式构造}
\label{sec:nlp-analysis-typed-logic}

令$e,t,\mathbb N$分别表示实体、真假和自然数类型。$G,R:e\to t$，$\mathrm{help}:e\to e\to t$，$\mathrm{count}:(e\to t)\to\mathbb N$；否定只接受真假值，量词绑定实体变量。一个构造器可从“数值表达式”槽开始，展开为$\mathrm{count}(P)$，再将谓词槽$P$展开为$\lambda x.G(x)\land R(x)$。每次只允许填入类型匹配的节点，且变量必须已绑定。

“有几位来宾已报名”的动作依次建立COUNT、LAMBDA、AND、G、R及两次变量引用，得到
\[
 z_1=\mathrm{count}(\lambda x.G(x)\land R(x)).
\]
在三个实体上求值，谓词向量为$(1,1,0)$，计数为2。产生式得分来自“文本--目标树”监督，训练可以对参考动作做交叉熵，束搜索累加动作对数概率并屏蔽类型不匹配的展开。例如COUNT后不能填实体常量a，因为a不是一元谓词；而$\mathrm{count}(\lambda x.G(x))$类型正确却遗漏报名条件，得到3。

否定的树位置改变意义。“并非所有来宾都已报名”为$\neg\forall x(G(x)\Rightarrow R(x))$，在当前域上为真；“所有来宾均未报名”为$\forall x(G(x)\Rightarrow\neg R(x))$，则为假。二者都能通过类型检查，只有保留否定与量词作用域才能区分。再比较“每位来宾帮助某位来宾”中的$\forall x\exists y$与“有一位来宾受到所有人帮助”中的$\exists y\forall x$：若帮助关系改为$\{(a,b),(b,c),(c,a)\}$，前者真而后者假。候选文法必须允许表达这种区别，才可能用语义监督学会正确选择。

\subsubsection{CCG与类型化\texorpdfstring{$\lambda$}{lambda}表达式}
\label{sec:nlp-analysis-ccg}

类型化树还可以与词语的句法组合一起产生。\term{组合范畴语法}（combinatory categorial grammar, CCG）给词项同时分配句法范畴和语义函数；$A/B$表示右侧接收B后成为A，$A\backslash B$则接收左侧B。函数应用在组合范畴的同时计算语义，因而能够追踪最终关系来自哪个词项\scite{nlp-analysis-ccg2005}

对补充句“小王帮助小李”，指定两个名字为$\mathrm{NP}:a$与$\mathrm{NP}:b$，“帮助”为
\[
 (\mathrm S\backslash\mathrm{NP})/\mathrm{NP}
 :\lambda y.\lambda x.\mathrm{help}(x,y).
\]
先向右应用b，得到$\mathrm S\backslash\mathrm{NP}:
\lambda x.\mathrm{help}(x,b)$；再接左侧a，$\beta$归约得到$\mathrm{help}(a,b)$。在给定关系外延中该式为真；若错误词项把帮助者和被帮助者颠倒，范畴仍可组合，却得到为假的$\mathrm{help}(b,a)$。句法成功没有替代词汇语义核验。

函数组合规则也可推迟实际参数的接入：
\[
 A/B:f,\quad B/C:g
 \quad\Longrightarrow\quad A/C:\lambda z.f(g(z)).
\]
为看清它与函数应用的差别，考虑问句“哪些来宾帮助小李”。用$\mathrm{VP}=\mathrm S\backslash\mathrm{NP}$表示一元谓词短语，$\mathrm{S_q}$表示以实体集合为答案的问句。设“哪些来宾”已经形成$\mathrm{S_q}/\mathrm{VP}$，语义为$f=\lambda P.\lambda x.G(x)\land P(x)$；“帮助”为$\mathrm{VP}/\mathrm{NP}$，语义为$g=\lambda y.\lambda x.\mathrm{help}(x,y)$。两者都尚未接到“小李”，可以先按上述规则形成$\mathrm{S_q}/\mathrm{NP}$：
\[
 \lambda y.f(g(y))
 =\lambda y.\lambda x.G(x)\land\mathrm{help}(x,y).
\]
此时才对小李的常量b做函数应用，得到$\lambda x.G(x)\land\mathrm{help}(x,b)$，在给定关系外延中恢复答案集合$\{a,c\}$。另一条推导可以先对“帮助”应用b，再交给f；两条推导给出相同语义。相对地，“未”作用于完整谓词“报名”得到$\lambda x.\neg R(x)$只是函数应用，不是刚才的函数组合。量化名词短语需要更高阶词项及允许的组合规则，不同量词作用域不能只凭词序猜出。

词项可由人工词典给定，也可从句子--逻辑式对提出候选并学习权重。概率CCG对“逻辑式$z$、推导$T$”联合评分$p(z,T\mid x)\propto\exp(\bm w^\top\bm f(x,z,T))$；同一逻辑式若有多个推导，应通过$\sum_Tp(z,T\mid x)$累加，训练只给逻辑式时同样边缘化推导。跨度图表保留范畴与语义项，束裁剪控制候选，但可能剪掉正确组合。未知词项、类型不一致及根范畴不符应返回无分析，不凭空补入一个事实谓词。

\subsubsection{逻辑式序列生成与答案监督}
\label{sec:nlp-analysis-logic-generation}

也可以把$z_1$序列化为\code{count(and(guest,registered))}，其中文法明确and对两个一元谓词逐实体取合取。文本--逻辑式对提供每步符号监督；解码用束搜索生成候选，再通过解析器恢复抽象语法树。无约束搜索可能输出括号不配对的串，文法约束则维护尚待填充的槽位与类型，过滤不能完成合法树的扩展。

假设某前缀后候选\code{registered}概率$.6$、\code{)}概率$.3$、其他符号总概率$.1$，而当前槽需要谓词，右括号会被屏蔽。保留词元原对数分数或在允许集合重归一化是不同搜索定义，必须固定；仅过滤后仍需完成所有括号、结束符及变量作用域。恢复树后遍历实体域，本例产生集合$\{a,b\}$和数值2，输出还应保留所求值的树，便于检查是不是正确问题。

只给答案$y$时，训练目标可改为
\[
 -\log\sum_{z:\,\operatorname{eval}(z,D)=y}p_\theta(z\mid x),
\]
其中$D$为固定事实环境。该和通常仅在搜索找到的候选集上近似；如果正确候选从未进入搜索，答案监督不会自动构造它。设正确计数式概率$.4$、常量表达式2概率$.3$、错误计数3概率$.3$，观察答案2时前两项都获支持，损失为$-\log .7\approx.357$。常量2是“答案碰巧正确”的伪解。把报名集合改为$\{a\}$，参考答案变1，正确计数式仍成立而常量失败。反例环境、多个问题或逻辑式约束有助于排除伪解，单次答案相等不构成语义证明。

\subsubsection{SQLNet查询骨架与槽填充}
\label{sec:nlp-analysis-sqlnet}

SQL的某些部分没有固定顺序，例如AND连接的两个筛选条件互换通常不改变结果。SQLNet针对单表查询骨架预测SELECT聚合、选择列以及WHERE条件集合，避免把条件的任意排列当作不同正确答案。其论文范围内的条件由列、$\{=,<,>\}$运算符和值组成，不能直接推广为任意连接或嵌套SQL\scite{nlp-analysis-sqlnet2017}

在registrations表上先问“ML-0418有几条已确认报名记录”。骨架为\code{SELECT AGG(COL) FROM registrations WHERE CONDS}。模型编码问题与列名，用列表示对问题位置做注意力：
\[
 \alpha_{c,i}=
 \frac{\exp(\bm c^\top\bm W\bm h_i)}
      {\sum_j\exp(\bm c^\top\bm W\bm h_j)},\qquad
 \bm q_c=\sum_i\alpha_{c,i}\bm h_i.
\]
这样event列可以关注活动编号，status列关注“已确认”。条件列用多标签选择，并另预测条件数量；列确定后预测相应运算符与问题中的值序列，SELECT列与聚合头也分别受监督。论文的值预测使用指针式生成；不能只给列名却不说明条件值从哪里恢复。

设条件数为2，event、status、pid的入选概率为$.9,.8,.1$，取前两列；运算符均选等号，值分别从问题恢复ML-0418，并按本节给定状态词典把“已确认”规范化为confirmed。SELECT列为pid、聚合为COUNT，得到：
\begin{codeblock}[language=SQL]{骨架恢复的单表查询}
SELECT COUNT(pid) FROM registrations
WHERE event = 'ML-0418' AND status = 'confirmed';
\end{codeblock}
三行中只有第一行满足两个条件，结果为1。交换条件顺序仍得1，训练条件集合时不应惩罚这种排列。这里状态词典是显式元信息，不能宣称原始SQLNet会无监督知道所有中文状态编码。单表结果尚未检查是否校外来宾；要回答原问题需要访问people，已经超出这个骨架。

\subsubsection{RAT-SQL模式关联与语法树解码}
\label{sec:nlp-analysis-ratsql}

多表查询必须知道问题里的对象对应哪张表、哪一列，以及表之间怎样连接。RAT-SQL把问题词、表和列共同编码，为同表、主外键、名称匹配、值匹配等关系加入专门表示，再由树解码器选择语法产生式和模式元素。关系感知注意力仍允许所有输入元素交互，不是只沿外键边传递信息\scite{wang-etal-2020-rat}

一层关系注意力把普通键向量$\bm k_j$换成$\bm k_j+\bm r^K_{i,j}$，把值向量换成$\bm v_j+\bm r^V_{i,j}$。设问题词“校外”对affiliation与status的普通点积分数均为0，而已提供的值匹配关系给前者额外2分，二选一注意力成为$.881,.119$；这说明模式关联怎样影响表示，不是模型准确率。原始RAT-SQL主要使用问题与库值的词面匹配；本例中文词语到external、confirmed的映射由给定业务词典提供，作为补充的模式链接输入。

解码器以深度优先顺序展开SQL抽象语法树，动作为应用产生式、选择列、选择表等，训练最小化参考动作序列的负对数似然。对当前问题，树须包含COUNT、people与registrations两个表、外键等值连接，以及活动、状态、身份三项筛选。列指针指向模式中的唯一列身份而非仅输出“pid”字符串，再由确定性序列化恢复别名和SQL：
\begin{codeblock}[language=SQL]{连接与聚合的语义恢复}
SELECT COUNT(DISTINCT p.pid)
FROM people AS p
JOIN registrations AS r ON r.pid = p.pid
WHERE r.event = 'ML-0418'
  AND r.status = 'confirmed'
  AND p.affiliation = 'external';
\end{codeblock}
常量槽由前面检出的文本值和业务词典绑定；这一步是教学执行接口，不能把RAT-SQL的结构解码描述成已经包办所有值识别和方言细节。执行连接得到3条匹配记录，活动筛选留下2条，状态筛选留下1条，该人的affiliation为external，去重计数为1。

本例主键已经防止同活动重复报名，DISTINCT不会改变正确结果；如果模式允许多条状态历史，它就可能变得必要。遗漏外键导致笛卡尔积，漏掉活动条件则把小赵在OTHER中的确认状态误算进来，结果变2。两者都不是语法错误。应检查名称或值是否正确链接，再检查树中是否选择正确的连接与条件；训练库与测试库不同，还需防止依赖固定列顺序的捷径。

\subsubsection{PICARD增量解析约束}
\label{sec:nlp-analysis-picard}

直接用语言模型生成SQL时，语法错误会占据束中的宝贵位置。PICARD在解码的每一步对候选前缀做增量解析，给无法继续成为合法查询的扩展赋$-\infty$；它改变推断中的候选允许性，而不要求重新训练一个专用动作词表。原文区分词法检查、无guards的语法解析和带guards的作用域检查，后者还要求引用列所属的表最终进入相应FROM范围\scite{scholak-etal-2021-picard}

例如前缀为\code{SELECT COUNT(}，模型提出\code{FROM}、\code{*}、\code{pid}，概率为$.50,.30,.20$。当前解析状态需要聚合参数，FROM不能使此前缀合法，故拒绝；保留后两项的原有分数，束宽1时选择星号。随后生成右括号、FROM registrations及筛选条件，才能完成查询。模型词元可能只是SQL单词的一部分，解析器要区分“暂时未读完”和“已经不可能完成”，不能要求每个子词前缀都单独是一条完整SQL。

若生成\code{SELECT p.status FROM people AS p}，status虽然存在于另一张表，却不属于p指向的people；模式及作用域检查应拒绝这一组合。相反，\code{SELECT COUNT(*) FROM people}完全合法却回答了“所有人有几人”，不是原问题。原文未实现所有可想象的类型检查，不能把它称为对任意SQL类型及业务语义的完整验证器。

若正确词元不在受检的top-$k$中，过滤无法把它重新加入；全部候选被拒绝时，应扩大候选或返回搜索失败。生成后的验证只能丢弃成品，增量检查可以提前保留搜索容量，但两者都不验证文本意图。边界状态、结束符接受条件与别名作用域是实现中应直接测试的对象。

\subsubsection{执行引导查询解码}
\label{sec:nlp-analysis-execution-guided}

语法检查不能发现全部运行问题，执行引导把数据库反馈也接入候选筛选。对可以独立执行的子查询或完整查询，使用只读连接、固定数据快照、超时和行数限制试运行；遇到不存在的列、函数或不满足数据库要求的组合时排除候选。不能把任意未完成SQL前缀直接交给数据库，再把“尚未写完”的语法错误当作该候选永远不可行。

固定本节数据库，三个候选分别是：把COUNT误写成不存在函数的查询、正确连接但漏掉活动条件的查询、完整参考查询。第一项在指定引擎中报错，后两项可运行，结果分别为2和1。执行引导可排除第一项，却无法只凭“成功执行”判断后两项谁回答了原问题。非空结果约束也只有在任务明确保证答案非空时才合理；“有多少已确认的校内来宾”的正确答案在当前表中就是0。

偶然相同结果更隐蔽。若删除OTHER那条记录，遗漏活动条件的错误查询也会返回1；将OTHER记录恢复后才显出差别。因此需要检查查询结构及多个反例数据库，而不是只比较当前执行值。候选模型可以利用检索到的问题--SQL示例，但示例中的表列和常量必须重新链接，不能把其他数据库的片段直接复制执行。

多轮查询还改变当前问题的完整含义。上一轮是“已确认报名的校外来宾有几人”，下一轮“那校内呢”应承接活动ML-0418与confirmed条件，仅将affiliation改成internal，结果为0。若历史只问“活动何时举行”，同一句“那校内呢”则没有足够条件恢复计数查询，应请求澄清。历史重写可以形成显式完整问题，也可以由模型读取前序问题与SQL；两种接口都应保留承接来源，不能在状态中悄悄遗失筛选条件。

\section{篇章分析}
\label{sec:nlp-analysis-discourse}

把每句话分析正确，仍可能不知道后一句的“他”是谁，也不知道它是在解释、反驳还是补充前一句。\term{篇章分析}（discourse analysis）处理跨表达单位的联系。本节按联系的对象区分提及身份、篇章单元关系、论证关系和内容连贯，而不是按使用了多长的编码器划分任务。长文档截断、候选筛选和历史范围会影响这些任务，却不是新的关系类型。

图~\ref{fig:nlp-analysis-discourse}采用补充段落“小王借到一本书。他读完后，归还了这本书。虽然封面已经磨损，但内容仍然清晰”。“小王”与“他”指同一个人，“一本书”与“这本书”指同一本书；“封面”与书是部分与整体的联系；让步关系则连接两个陈述。三种边的端点和意义不同。

\begin{figure*}[htbp]
  \centering
  \begin{tikzpicture}[
  x=1mm,y=1mm,
  every node/.append style={font=\figurefont\footnotesize,align=center},
  mention/.style={draw=FigureInput,fill=FigureInput!10,
    rounded corners=1mm,inner sep=1.5mm,minimum height=8mm},
  identity/.style={draw=FigureInput,line width=1.1pt},
  bridge/.style={->,draw=FigureOutput,dashed,line width=1.1pt},
  discourse/.style={->,draw=FigureModel,dash dot,line width=1.1pt}
]
  \node[mention] (wang) at (20,0) {小王};
  \node at (42,0) {借到};
  \node[mention] (book) at (71,0) {一本书};
  \node at (86,0) {。};
  \node[mention] (he) at (20,-20) {他};
  \node at (42,-20) {读完后，};
  \node[circle,draw=FigureInput,line width=1.1pt,
    minimum size=3mm,inner sep=0] (zero) at (67,-20) {};
  \node at (90,-20) {归还了};
  \node[mention] (book2) at (128,-20) {这本书};
  \node at (145,-20) {。};

  \draw[identity] (wang.south) -- (he.north);
  \draw[identity] (wang.south) -- (20,-10) -- (67,-10) -- (zero.north);
  \draw[identity] (book.north) -- (71,9) -- (128,9) -- (book2.north);
  \node[anchor=south] at (113,1) {同一本书};
  \node[anchor=north] at (67,-23) {省略位置};

  \draw[draw=FigureModel,fill=FigureModel!5,rounded corners=1mm]
    (2,-59) rectangle (84,-43);
  \draw[draw=FigureModel,fill=FigureModel!5,rounded corners=1mm]
    (94,-59) rectangle (162,-43);
  \node at (12,-51) {虽然};
  \node[mention,draw=FigureOutput,fill=FigureOutput!10] (cover) at (32,-51) {封面};
  \node at (62,-51) {已经磨损，};
  \node at (128,-51) {但内容仍然清晰。};
  \draw[bridge] (cover.north) -- (32,-35) -- (128,-35) -- (book2.south);
  \node[fill=white,inner sep=1mm] at (102,-35) {PART-OF};
  \draw[discourse] (43,-59) -- (43,-69) -- (128,-69) -- (128,-59);
  \node[fill=white,inner sep=1mm] at (85,-69) {让步：两个陈述};

  \draw[identity] (3,-82) -- (15,-82);
  \node[anchor=west] at (18,-82) {同一身份};
  \draw[bridge] (55,-82) -- (67,-82);
  \node[anchor=west] at (70,-82) {桥接};
  \draw[discourse] (109,-82) -- (121,-82);
  \node[anchor=west] at (124,-82) {篇章关系};
\end{tikzpicture}

  \caption{同一段落中的身份共指、桥接与篇章关系。实线连接同一对象的提及，虚线连接“封面”与书的部分--整体关系，点划线连接篇章单元。空心位置标记不是原文字符，而是“归还”前检测出的省略主语位置，先行项为小王。共指可以形成等价簇，桥接和篇章关系不能并入该簇。}
  \label{fig:nlp-analysis-discourse}
\end{figure*}
\FloatBarrier

\subsection{提及、先行项与共指}
\label{sec:nlp-analysis-coreference}

\term{提及}（mention）是文本中指向对象的一次表达，其位置可以是“机器学习读书会”这样的名词短语，也可以是代词。\term{共指}（coreference）连接指向同一对象的提及；\term{先行项}（antecedent）是用于解释当前表达的较早表达。共指簇描述身份，先行项指针描述一种恢复路径，二者不是同一种数据结构。

“一本书”的封面不是那本书本身，因此“封面”到书应输出PART-OF桥接边，不进行身份合并。“活动改期了。这让小王调整了行程”中的“这”指向改期这一事件或前述内容，是话语指示，不一定存在一个对应的实体名词跨度。知识库链接又要求选择库中身份，见第~\ref{chap:nlp-information-extraction}章；文档内已经共指，不等于库中一定有对应项。

中文还可能没有显式代词。图中“归还”前应先检测一个省略位置，再为该位置选择小王。可在字符边界上训练“是否省略论元”的分类器，以谓词、句法和上下文特征形成候选，位置用锚点而非伪造字符跨度表示。若省略位置是人工给定的，所得先行项准确率只是条件结果；端到端评价还要计入位置漏检、误检与身份错误。

一种基线对提及对训练二分类，再按阈值凝聚聚类。若A--B与B--C的概率为$.9,.8$，A--C只有$.1$，阈值$.5$并取连通分量仍会把三者全合并；传递闭包修复了格式，却可能放大一次错连。联合簇预测可以把这些冲突放入目标，生成式方法则必须将生成的提及名单重新对齐原文、去重并恢复簇，不能将两个同名跨度混为一处。跨文档时先按名字、时间和来源阻塞候选，再做身份匹配；阻塞漏掉的别名不可能被后续分类补回，同名也不足以证明同一人。

\subsubsection{Hobbs句法先行项搜索}
\label{sec:nlp-analysis-hobbs}

Hobbs的朴素算法不是简单选择最近名词，而是按确定顺序遍历句法树，再用数、性等一致性条件筛选候选。它从直接支配代词的NP开始，上行到第一个NP或S，记为X，并保存上行路径p；在X下、p左侧按从左到右的广度优先顺序访问，只提出与X之间还隔着NP或S的名词短语。这一附加条件用于排除同一局部结构中的若干不合法候选\scite{nlp-analysis-hobbs1978}

若尚未到全句最高S，就继续上行到下一个NP或S。新X是NP且路径没有经过它直接支配的$\bar N$时，提出X本身；再广度优先搜索路径左侧，提出遇到的NP。若X为S，还搜索路径右侧，但遇到NP或S后不深入其内部。重复上行；到达最高S仍未找到时，按由近及远的顺序遍历前句，每棵树内部仍从左到右广度优先提出NP。候选逐个经过一致性及已知非共指约束，首个接受者成为结果，候选用尽则返回未解析。

以“She thanked him”为当前句、上一句“Mary met John”为例，假设只考虑这两句，且已知Mary使用she、John使用he。当前句从him的NP上升到S；左侧She的NP与S之间没有另一个NP或S，因此第一轮不提出She。当前S已是最高S，转向上一句；广度优先先提出主语Mary，但性特征不匹配him，随后提出宾语John并接受。输出是him到John原文提及的指针，不是句法树节点编号。

这个例子不需要模型训练，信息来源是句法树、规则和一致性词汇。若句法树错把短语嵌套关系改变，候选次序也随之改变；如果两个候选都满足形式条件，语义和常识仍可能推翻首选。英语的性、数及$\bar N$结构不能原样移植到所有语言。中文零代词更没有可直接定位的代词NP，必须先检测省略位置并建立相应句法输入；桥接或话语指示也不应通过放宽筛选就混入身份簇。

\subsubsection{端到端跨度共指与先行项排序}
\label{sec:nlp-analysis-span-coreference}

跨度共指把提及检测与连接一起评分。Lee等的方法枚举有最大宽度限制的跨度，以两端上下文、内部注意力和宽度构造$\bm g_i$；提及分数$s_m(i)$用于裁剪，对保留提及只考虑较早候选及虚拟先行项$\epsilon$。两真实提及的分数为
\[
 s(i,j)=s_m(i)+s_m(j)+s_a(i,j),\qquad s(i,\epsilon)=0,
\]
其中$s_a$利用两跨度表示、交互特征、距离与可用元信息。虚拟项表示开始新簇，或该候选不是有效提及；如何输出单例还要遵循评测规范\scite{lee2017coreference}

对候选集合$Y(i)$做softmax。一个提及可能有多个正确先行项，若$G(i)$为当前可见候选中与它同簇的先行项集合，训练损失是
\begin{equation}
 -\sum_i\log\sum_{j\in G(i)}p(j\mid i,x).
 \label{eq:nlp-analysis-coreference-marginal}
\end{equation}
没有可用的正确先行项时取$G(i)=\{\epsilon\}$，包括首提及和裁剪掉全部正确先行项的情况；不能无条件要求模型连接到不存在的候选。若把多个正确先行项任意挑成唯一标签，会额外惩罚同簇的其他合理连接。

设文本已明确“小王”是“王老师”的别称，依次出现$m_1=\text{王老师}$、$m_2=\text{小王}$、$m_3=\text{他}$。给$m_2$到$m_1$的分数2、到$\epsilon$的分数0，选择$m_1$。对$m_3$，到$m_1,m_2,\epsilon$的分数为$2,3,0$，对应概率约为$.259,.705,.035$；两个真实先行项都正确，训练使用合计约$.965$，推断选$m_2$。把指针看作无向连接取连通分量，恢复$\{m_1,m_2,m_3\}$，与直接选$m_1$的簇相同。

下一次出现“另一位王老师”时，若到前三者的分数为$-2,-3,-1$，虚拟项0胜出，建立不同身份。仅比较名字会错误合并；正确词义和正确人名跨度也阻止不了这种身份错误。粗到细候选筛选可先用廉价内积缩小先行项集，再使用完整配对评分，降低计算量但提高遗漏风险。应检查提及召回、簇内误合并与跨簇漏连接；只报指针准确率不能完整反映共指簇质量。

\subsection{篇章单元与关系}
\label{sec:nlp-analysis-discourse-relations}

身份共指连接对象，篇章关系连接承载内容的单元。\term{基本篇章单元}（elementary discourse unit, EDU）是标注规范选定的最小分析片段，常接近分句，但不等于“按每个逗号切开”。PDTB式分析标注连接词、两个论元及关系，不要求全文形成同一棵树；RST式分析还要求核心／附属地位及覆盖全文的树。前者的论元也可能跨越多个基本单元。

用“下雨了，所以路面湿了。车辆减速。”作为统一材料，本节教学分割为三个单元：$u_1$“下雨了，”、$u_2$“所以路面湿了。”、$u_3$“车辆减速。”。前两者有显式连接词，后两者可以标为隐式结果关系；此处是对文字表达关系的解释，并不从“所以”证明道路实际因降雨变湿。DISRPT把分割、连接词检测、关系分类分别建成任务，这些子任务分数不能充当完整篇章树的分数；完整来源见文献说明。

\subsubsection{篇章边界标注与论元对关系分类}
\label{sec:nlp-analysis-discourse-boundary}

先按参考EDU边界为字符或词生成标签，训练逐位置分类器或CRF。若使用B/I分别表示单元开始和内部，以上16字符序列在位置0、4、11取B，其余取I；句首强制B，解码后恢复$[0,4)$、$[4,11)$、$[11,16)$。标点在本例归入其左侧单元，其他规范可以不同，但原文覆盖须无遗漏且无重叠。

连接词检测从词典提出“所以”$[4,6)$，再根据上下文判断是否确实起连接作用；仅命中词典不足以消解同形的非连接用法。随后提出论元对，编码$u_1,u_2$及允许上下文，对CAUSE、CONTRAST、ELABORATION等关系做softmax交叉熵训练。若三个logit为$2,0,-1$，CAUSE概率约$.844$，恢复“原因论元$u_1$、结果论元$u_2$、连接词所以”的记录。标签名和方向需匹配实际资源。

对$u_2,u_3$，连接词字段为空，假设CAUSE与无关系得分为$1.5,0$，前者概率约$.818$，输出隐式结果关系。本例没有给出真实因果验证，只说明模型认为这种承接最合适。若选错论元为$u_1,u_3$，即便CAUSE标签正确也不是同一个输出；应分开评价边界、连接词、论元跨度与关系。关系训练若使用人工论元，测试也须说明是否同样给定，不能隐藏论元恢复的难度。

\subsubsection{RST移进--归约篇章解析}
\label{sec:nlp-analysis-rst-transition}

RST需要把全部单元组合成层次树。这里采用二叉教学规范：叶节点为给定EDU，内部节点记录关系及左右孩子的核性；N为核心，S为附属，允许NS、SN或相应多核心关系的NN。核心是当前表达组织中的中心内容，不表示事实更真。多叉参考树必须按固定规则二叉化，评测时再按规范恢复。

状态由子树栈和未读EDU队列组成。SHIFT将下一个单元入栈，REDUCE$(r,\nu)$弹出相邻的两个子树，按关系$r$和核性$\nu$组成父树；队列为空、栈中恰好一棵覆盖全文的树时结束。参考树经后序遍历给出动作监督，分类器根据栈顶、下一单元、边界和全文表示产生合法动作logit，训练以参考动作交叉熵为目标。栈不足两个不能归约，跨越非相邻区间不能任意拼接。

对三个单元，先执行SHIFT$(u_1)$和SHIFT$(u_2)$，再以REDUCE(Cause,SN)构成$[u_1^{S},u_2^{N}]$。随后执行SHIFT$(u_3)$，再次以REDUCE(Cause,SN)将前面的子树作为原因背景，与核心结果$u_3$组成整树。若三次SHIFT各得1分、两次归约得2与3分，总分8；另一条先合并$u_2,u_3$再接$u_1$的轨迹，若归约分数为2与1，总分6。束搜索比较完整动作路径，贪心则可能提前作出不可逆组合。

还原树后，检查叶顺序恰为$u_1,u_2,u_3$，每个单元只出现一次，每个内部节点关系及核性均合法。单元分错可能让所有上层跨度改变，关系名预测错则保留树形却改变解释，二者应分别评价。局部分类全部有输出不代表完成了树恢复；剩两个不连通子树时不能悄悄称为完整RST解析。

\subsubsection{递归分割式篇章树恢复}
\label{sec:nlp-analysis-rst-split}

也可以从整段跨度向下决定划分。令$[i,j)$为EDU边界区间；若$j-i=1$，直接返回该叶节点，否则在$i<k<j$上预测分割位置，再预测核性$\nu$及关系$r$，递归处理$[i,k)$与$[k,j)$。编码器可汇聚整段及两侧上下文，分割头对候选边界做softmax，核性与关系头在选定分割条件下训练。参考树提供各内部节点的三项监督。

对$[0,3)$，设分割$k=1,2$的logit为1、3，贪心选$k=2$，概率约$.881$；核性SN与关系Cause得分最高。对子区间$[0,2)$，唯一分割$k=1$，仍取Cause/SN；$[2,3)$及两个更小区间都是叶节点，停止。恢复的树与上一小小节相同。该例实际选择两次分割，三个单元全部被覆盖。

若要对可加、无历史依赖的局部分割分数求全局最优，可以计算
\[
 D(i,j)=\max_{i<k<j,\,r,\nu}
       \{s(i,j,k,r,\nu)+D(i,k)+D(k,j)\},
 \quad D(i,i+1)=0,
\]
并保存分割及标签回溯。普通自顶向下贪心并不自动执行这一全局动态规划；解码器若依赖此前分割历史，上式也不能直接复用。高层若误选$k=1$，其子区间再怎么细分都不能出现正确的$[0,2)$成分，这说明早期分割错误为何会影响整树。核性与关系指标仍按RST规范，不能把共指或桥接硬塞进这棵树。

\subsection{论点、理由与论证结构}
\label{sec:nlp-analysis-argumentation}

篇章的“解释”关系不等于一个理由已经足以支持结论。\term{论证分析}先定位观点与理由，再恢复支持、反驳等有方向的联系。以下补充材料含三个单元：$a_1$“应增加自习座位”，$a_2$“晚间常有人找不到空位”，$a_3$“本月晚间到场人数持续上升”。在本节教学解释中，$a_2$支持$a_1$，$a_3$为$a_2$提供背景依据；这不证明统计为真，也不证明增加座位是唯一解决办法。

原文跨度和单元类型可由序列标注得到，关系分类以单元对为输入。流水线可用SVM或编码器分类器；生成式输出也必须还原到原文位置，明确“理由指向结论”的方向。隐含结论没有显式跨度，跨段依据不能被当前段落窗口覆盖，都需要单独标注。以下结构约束只适用于明确规定的树或森林，不代表所有自然论证都如此。

\subsubsection{论证单元分类与ILP结构恢复}
\label{sec:nlp-analysis-argument-ilp}

Stab与Gurevych先识别单元、分类论点与前提、预测连接，再用整数线性规划协调局部输出；边权结合关系预测、由入出边导出的论点倾向及类型预测。恢复结构后，无出边者作为论点，其余作为前提，支持／反驳立场另行分类。其段落层约束允许一棵或多棵树，不应笼统宣称每个段落都强制连通；完整来源见文献说明。

用二元变量$x_{i,j}$表示选择从理由$i$到目标$j$的边，目标是最大化所选边的总权重：
\[
 \max\sum_{i\ne j}w_{i,j}x_{i,j}.
\]
例如将类型概率$p_j(\mathrm{Claim})$和关系概率$q_{i,j}$组合成边权，再在验证集调节相对权重；这是一种教学评分实现，不是把概率之和误称为联合概率。结构约束为
\[
 x_{i,i}=0,\quad \sum_jx_{i,j}\leq1,\quad
 \sum_{i,j}x_{i,j}\leq n-1.
\]
这三项还没有排除环。可另引入$0\leq u_i\leq n-1$的层次变量，要求
\[
 u_i\geq u_j+1-n(1-x_{i,j}),\qquad i\ne j.
\]
选边时层次严格下降，沿环累加会推出$u_i\geq u_i+k$的矛盾，故不能成环；未选边时约束在给定范围内松弛。这样得到有向森林。只有标注要求整段一棵树时，才增加恰好$n-1$条边及适当根条件。

对三个单元，禁止$a_1$出边，其他有限候选为$w_{2,3}=6,w_{3,2}=5,w_{2,1}=4,w_{3,1}=2$。局部选择$a_2\to a_3$与$a_3\to a_2$得11分却成环；合法组合$a_2\to a_1,a_3\to a_2$得9分，另一组合$a_2\to a_3,a_3\to a_1$得8分，最优取9。恢复边后，$a_1$为根论点，$a_2,a_3$为前提，再给边附上支持标签和各单元原文跨度。

这一步解决的是结构冲突，不是支持强度验证。两个独立理由共同支持同一结论时可以各有一条入向该结论的边；若只有联合使用多个前提才成立，普通二元边不能充分表达该合取，需要额外论证节点或超边。把同一单元既作子论点又作前提的情况压成一个全局类型，也可能损失信息，须与标注层级一致。

\subsubsection{双仿射论证依存解析}
\label{sec:nlp-analysis-argument-biaffine}

给定论证单元时，可将每个单元编码为向量，为“依据端”和“目标端”分别投影，再复用第~\ref{sec:nlp-analysis-biaffine}节的双仿射评分；类型头预测论点、前提、非论证，关系头预测支持、反驳或无边。训练对标注连接及类型计算损失，解码再施加本节选择的森林约束。这里的节点是论证单元，标签不是句法nsubj、obj。

令二维依据向量$\bm h_2=(1,0)$、$\bm h_3=(0,1)$，目标向量$\bm g_1=(4,2)$、$\bm g_2=(0,5)$、$\bm g_3=(6,0)$，双线性矩阵取单位矩阵、其他项为0，即得到上一例四个边分数$4,2,5,6$。用同样的无环约束恢复9分结构，再用各选中边的关系头选支持；换评分网络没有改变输出对象或取消解码责任。

若还需从原始词语恢复单元，可采用词元级结构编码：在同一单元内部用带类型的APP边连接相邻词，以特定端点代表整个单元，再在端点间预测论证连接。Ye与Teufel的代表表示用单元末词作为对外端点，并允许非树的论证图；其依存图箭头约定与“理由指向结论”的论证约定需要显式转换\scite{ye-teufel-2021-end}

例如词序“晚间／常／无／空位”形成内部链后，恢复连续跨度及Premise类型；其末词与“增加／座位”单元末词之间的预测连接，再映射为两个完整单元的支持边。缺一条内部APP边会拆坏单元，端点连错则改变支持对象。多父节点允许一个依据关联多个目标，但仍不能自动表达合取前提；隐含结论也需要独立节点约定。若规范允许图，不能再强行套单父树解码后宣称完整覆盖。

\subsection{内容衔接与连贯性}
\label{sec:nlp-analysis-coherence}

实体与关系可以作为分析输出，也可以为段落排序或连贯性判断提供依据。这里的输出是一个分数或句子排列，不是共指簇。用原文和受控扰动训练时，要明确究竟交换句子、删除关键指代，还是插入无关材料；学会识别一种扰动，不等于能够判断全部写作质量。

比较三句材料：A“小王借了一本书”，B“他读完了这本书”，C“这本书被归还”。原顺序为ABC。交换成ACB会让归还先于阅读，却仍保留几乎全部实体词；删去B里的“他”会破坏显式身份衔接，插入无关句又是不同故障。编码器可直接分类原文与扰动文本，但要排除句首编号、段落长度及原始位置等捷径。更长上下文或分层编码改变可见信息与成本，不改变验收对象。

\subsubsection{实体网格与排序学习}
\label{sec:nlp-analysis-entity-grid}

\term{实体网格}以句子为行、共指实体为列，用S、O、X、--表示主语、宾语、其他出现和缺席。同一句多次出现同一实体时，按S优先于O、O优先于X的规则选择一个格值。Barzilay与Lapata用相邻句的角色转移分布作为特征，再学习原文优于扰动顺序的排序函数\scite{barzilay-lapata-2008-modeling}

对三句材料，将“他”合并到小王、“这本书”合并到书，网格为
\[
 \begin{array}{c|cc}
    &\text{小王}&\text{书}\\ \hline
 A&\mathrm S&\mathrm O\\
 B&\mathrm S&\mathrm O\\
 C&-&\mathrm S
 \end{array}.
\]
C是被动句，书的语义角色是受事，但句法角色为主语S。两列各有两次相邻转移，共4次，SS、S--、OO、OS各占$1/4$。设教学权重$w_{\mathrm{SS}}=2,w_{\mathrm{OO}}=1$、其余0，原顺序分数$.75$；顺序ACB产生S--、--S、OS、SO，分数0。权重在真实系统中由训练文本及扰动样本学习，不由本例的直觉固定为普遍规则。

可用排序SVM的损失$\max(0,1-f(D)+f(D'))$，其中$f(D)=\bm w^\top\bm\phi(D)$，使原文相对扰动文本取得间隔。当前两者分差$.75$，损失为$.25$；优化会调整权重，但仍需正则化和跨文档验证。如果漏掉“他”与小王的共指，网格变三列，SS消失，OO只占$1/6$，同一权重下原文分数降至$1/6$。共指错误因此会通过特征影响排序。

这种表示只捕捉局部实体延续。即使句子全围绕同一个实体，也可能重复、矛盾或缺少论证，网格未必区分。段落级应用还要固定行的单位，不能训练时一行一句、测试时一行一段却沿用同一转移解释。

\subsubsection{句对衔接评分与束搜索排序}
\label{sec:nlp-analysis-pair-ordering}

另一种方法为“哪句适合接在哪句后面”训练有方向的评分$s(a,b)$。正样本来自原文相邻句，负样本从非相邻或扰动位置抽取；损失可用二分类或在候选后继上的softmax。起始分数$s_0(a)$表示哪句适合开头，对排列$\pi$的目标为
\[
 S(\pi)=s_0(\pi_1)+\sum_{t=2}^m s(\pi_{t-1},\pi_t).
\]
没有重复使用任何句子是搜索约束，内容关系是否合理仍取决于评分。

取$s_0(A)=2,s_0(B)=0,s_0(C)=-1$，边分数为$s(A,B)=3,s(A,C)=4,s(B,C)=2,s(B,A)=-1,s(C,A)=-3,s(C,B)=-2$。束宽2时首步保留A的2分、B的0分；扩展后AC得6、AB得5、BC得2、BA得$-1$，保留AC与AB。第三步只能接未用句子，ACB得$6-2=4$，ABC得$5+2=7$，最终恢复ABC。句子ID再映射到完整原文，不能只拼接模型内部编码片段。

束宽1会在第二步选择AC，最终得到较差的ACB；可见局部最高分并不保证整体最佳。这个例子穷举6个排列可确认7为全局最大，真实长文档通常做不到全部枚举。每次扩展必须更新已用集合，遇到同文本的两条句子也按不同输入ID处理。标注多个合理顺序时，单一参考排列还可能惩罚可接受输出，应按内容关系另行核验。

\subsubsection{指针式句子排序}
\label{sec:nlp-analysis-pointer-ordering}

指针网络在每一步直接从剩余输入句子中选择一个。编码器得到全部句子表示$\bm h_j$，解码器将已选句子历史压入$\bm z_t$，打分$g_{t,j}=\bm v^\top\tanh(\bm W\bm h_j+\bm U\bm z_t)$。给已用句子的分数设为$-\infty$，对其余项softmax，训练用参考顺序每一步的交叉熵。与固定句对分数相比，当前选择可以依赖全部已选历史，而非仅最后一句。

若首步对A、B、C的概率为$.6,.3,.1$，选择A；第二步将A屏蔽，对B、C的概率为$.75,.25$，选择B；第三步只剩C，概率1。完整ABC路径概率为$.45$，对应训练损失$-\log .45\approx.799$。ACB路径概率为$.6\times.25=.15$。若用束搜索，保留的是不同历史及其剩余集合，不能把各步在不相同历史下的最大概率拼成一条路径。

当全部$m$个ID恰好选过一次时结束；提前结束、重复选择和漏句分别是格式错误。训练时给参考历史而测试时使用自身预测会产生历史分布差异，错误首句可能影响后续所有分数。长文档还受句数、全段编码与候选注意力成本限制。最终评价应同时看排列和事件、指代、论证关系，不把“输出了一个合法排列”当成连贯性已经解决。

\section{语用分析}
\label{sec:nlp-analysis-pragmatics}

“明天能完成吗”在制定计划时可以询问可行性，在对方多次延误后也可能带有催促。字面命题相近，交流中的作用却不同。\term{语用分析}（pragmatic analysis）利用说话者、被回应对象、历史话语与当前情境，判断表达承担什么作用，以及哪些未明说的内容得到背景支持。背景可以缩小解释范围，但不能为模型任意补写意图提供许可。

\subsection{话语行为}
\label{sec:nlp-analysis-dialogue-acts}

\term{话语行为}（dialogue act）用请求、告知、确认、承诺等标签表示一次表达在交流中的功能。疑问句不必只求信息，陈述句也可能构成请求。标签规范须说明是否允许多功能、是否按完整轮次或轮内片段标注，以及当时有哪些背景可供判断。识别当前用户行为与选择系统下一步行为，是两个需要分别验收的任务。

单轮基线可以把词语、句式、说话者及被回应对象等特征送入最大熵或SVM分类器。面对新标签，也可以把标签定义与例子作为候选描述，计算“当前话语--标签描述”的匹配分数；这改变的是类别信息来源，不表示模型无需理解标签边界。例如定义将“询问可行性”和“要求加快进度”明确分开后，才能比较两种历史对同一句话的作用。如果标注者看了后续回答而系统预测时只看到当前轮，训练目标已包含不可用信息，应重新限制标注依据或明确改为事后分析。

\subsubsection{跨轮HMM与CRF话语行为标注}
\label{sec:nlp-analysis-dialogue-sequence}

把每轮行为$a_t$作为状态，当前文字$u_t$作为观测，HMM使用$p(a_t\mid a_{t-1})p(u_t\mid a_t)$。转移概率由已标注相邻行为计数估计，观测概率可由行为条件的词语模型或离散表达模式估计。若只把分类器输出的$p(a_t\mid u_t)$直接当作发射概率而不作处理，就改变了概率模型的含义。

固定“明天能完成吗”这一观测，教学发射概率为$p(u\mid I)=.6,p(u\mid U)=.4$，I为询问、U为催促。已知前轮是计划讨论P时，转移到I、U的概率为$.8,.2$，两个乘积为$.48,.08$，归一化后I占$6/7$；已知前轮是延误解释E时，转移概率为$.2,.8$，乘积为$.12,.32$，U占$8/11$。同一句式的观测贡献不变，行为上下文改变排序。数值仅用于说明机制，不能把“出现解释”硬编码成必然催促。

前轮状态也不确定时，在线过滤先计算
\[
 p(a_t\mid u_{\leq t})\propto
 p(u_t\mid a_t)\sum_{a_{t-1}}
 p(a_t\mid a_{t-1})p(a_{t-1}\mid u_{<t}).
\]
若需要前缀的最高分整条行为序列，则把求和换成Viterbi的最大化并保存指针；最高分路径末标签与末状态边缘概率最大者并不总相同。这里处理的是已发生的轮次，不读取未来回答。

CRF则直接定义前缀条件分数，将文字、说话者及历史特征放进位置得分，并学习行为转移。对同一问句设$I,U$位置分数为$1,.5$，P后的转移分数为$1,0$，两候选总分$2,.5$，偏向I；E后的转移分数为$-.5,1$，总分$.5,1.5$，偏向U。其训练和回溯复用第~\ref{sec:nlp-analysis-bmes-crf}节，但话语行为一般没有BMES那样的固定边界掩码。

在线使用时，编码器和训练样本必须限制到当前前缀，不能使用整段双向编码得到的未来信息。事后标注完整对话可以利用全部轮次，但那是另一输入条件；相同的HMM或CRF名称并不保证二者可直接比较。话语行为序列模型的经典统计设置见Stolcke等的工作，其观察单位为已分割话语，不应省略分割条件后宣称直接处理任意原始对话\scite{stolcke-etal-2000-dialogue}

\subsubsection{层次对话编码与行为预测}
\label{sec:nlp-analysis-dialogue-hierarchical}

相同的行为标签历史可能包含不同事实，仅用行为转移不一定够。层次编码先在轮内汇聚词语，再在轮间汇聚已发生的历史。设$\bm h_t=f(u_t)$为当前文字表示，$\bm c_t=g(\bm h_{\leq t},\text{说话者、回应关系})$为因果历史表示，预测头在$[\bm h_t;\bm c_t]$上做softmax，使用当前轮行为标签训练交叉熵。轮内编码可双向读取当前已结束的文字；轮间编码必须屏蔽未来轮次。

用一个可复算的二维头说明背景作用。固定问句的基底得分为1，取历史标量$c$，询问与催促的logit分别为$1+c$、$1-c$。计划历史编码为$c=.8$时，分数为$1.8,.2$，询问概率约$.832$；延误历史编码为$c=-.8$时，催促概率约$.832$。不看历史的基线$c=0$给两类各$.5$。真实网络通过监督学习这些表示和参数，本例只暴露各信息怎样影响最终得分，不声称两种历史天然对应某个固定向量。

多方对话中，“明天能完成吗”可能回复较早的一位参与者，而不是紧邻上一轮；模型应输入允许的回应边或同时预测它。只有说话者ID、没有被回应对象时，长历史也可能混淆任务。“谢谢，能再发一份吗”同时致谢与请求，若规范允许两种功能，应使用多标签头或分段后分别恢复标签；把它压为单标签的歧义不能靠扩大模型消除。评测应分别检查行为标签和它对下一步回应选择的实际帮助。

\subsection{隐含意义与省略恢复}
\label{sec:nlp-analysis-implicit-meaning}

话语还可能省略了需要询问的对象，或借字面表达暗示另一种要求。“那明天呢”没有独立说明问的是开放时间、天气还是能否参加。此时应从允许背景形成候选解释，再选择、保留歧义或请求澄清，不能把语言模型最容易补出的句子当作已获支持的意图。

构造两段独立教学历史。历史H1的前问是“今天能到图书馆借书吗”，H2的前问是“今天的空气质量如何”；当前表达均为“那明天呢”。根据平行结构可分别提出“明天能到图书馆借书吗”和“明天的空气质量如何”，另设背景不足候选。用$(H,u,\text{候选解释})$句对或多段输入训练分类、排序头，监督来自标注的合理补全；若H1下三类logit为$2,0,-1$，选第一项，概率约$.844$，H2下前两项对调则选第二项。没有历史时，两种补全均缺少依据，应输出未知，而不是只在两者中强行择优。

候选也可从对话槽位恢复。H1保留动作“借书”、地点“图书馆”、疑问类型“是否可行”，当前轮只替换日期槽；H2保留主题“空气质量”和属性询问。“小王借了书，小李也借了”则可从平行句法提出省略宾语“书”，但仍要区分同类书与同一本书，词语补全不能自动完成身份共指。零代词的位置与身份分析回用第~\ref{sec:nlp-analysis-coreference}节。

带复制机制的序列到序列补全可以从历史复制保留片段，再生成变化部分。每步输出分布为
\[
 p(w)=\lambda p_{\mathrm{vocab}}(w)
 +(1-\lambda)\sum_{i:x_i=w}\alpha_i.
\]
设当前要恢复“图书馆”，词表概率$.2$、历史中该词的复制注意力总和$.8$、$\lambda=.25$，混合概率为$.65$。复制路径说明内容来自历史，但只复制并不证明角色正确；把“图书馆”复制到人物槽仍是错误。训练用完整补全监督，恢复后检查日期替换、保留的谓词与论元，以及新增词是否由历史或任务词典支持。

上下文查询重写把上述恢复结果交给检索或SQL模块，受约束解释生成则可要求输出“补全内容＋来源片段＋未确定条件”。沿全章通知，咨询发生在2025年4月17日，“我明天下午到可以吗”中的明天可规范化到4月18日，下午却不能直接改写成精确的14:00。对于校外来宾，还缺少是否已在4月16日18:00前提交报名信息的状态，应保留这个条件或询问，而不是生成无条件的参加许可。

隐含意图又不同于省略槽位。“这里有点冷”在天气讨论中可能只是报告感受，在临近开窗的人时可能暗示关窗；相同字面内容允许不同解释，后者也不由逻辑蕴含保证。自然语言推断模型可以筛查候选是否与背景冲突，却不能把“背景不反驳”等同于“说话者确有此意”。合理候选并存时保留多种解释及分歧，只有额外背景足够时才收窄。输出位置、补全内容、背景支持和回应是否得当，需要分别核验。

\section*{文献说明}

UD的跨语言标注框架与图方法的依存树解码分别依据以下资料。两者提供标注对象及结构恢复方法，具体的单根约束、得分和教学算例仍以本章声明为准。

\begin{fullwidth}
\raggedright
\fullcite{nivre2020universal}\par
\fullcite{mcdonald-etal-2005-non}
\end{fullwidth}

篇章任务的划分与论证森林的联合恢复分别参考以下工作。本文给出的数值与单元划分均为教学构造。

\begin{fullwidth}
\raggedright
\fullcite{braud-etal-2023-disrpt}\par
\fullcite{stab-gurevych-2017-parsing}
\end{fullwidth}

数据库任务的输入条件与数据划分可查阅以下原始资料。这里集中列出完整书目，避免长作者列表挤占正文边注；前文仅借它们说明任务范围，没有将不同基准的结果混作同一性能比较。

\begin{fullwidth}
\small
\fullcite{yu-etal-2018-spider}\par\medskip
\fullcite{nlp-analysis-spider2025}\par\medskip
\fullcite{yu-etal-2019-sparc}\par\medskip
\fullcite{yu-etal-2019-cosql}
\end{fullwidth}

\section{本章小结}
\label{sec:nlp-analysis-summary}

本章开头的问题是：圈出通知中的时间和人物之后，还缺少哪些可以恢复、检查并传给后续任务的结果。词法给出单位边界和局部属性，句法给出组合或依存结构；词义、角色、概念图及形式表达分别回答哪种用法、怎样参与、概念如何联系以及如何求值。篇章把跨句身份与内容联系起来，语用再根据可见背景约束话语功能和省略解释。它们互相提供依据，却不能互相冒充交付结果。

Viterbi和CKY说明了一个共同机制：将局部评分放进明确的合法结构空间，保存最优子问题和回溯信息，才能恢复整句结果。图解析、共指聚类、查询构造和句子排序继续使用这一思想，但合法性随对象而变。树的单中心词约束不适用于共享论元图，身份簇的传递性不适用于桥接，查询可执行也不表示符合原问题。

实际系统应同时保留原文位置、标注规范、可用背景及恢复轨迹，让候选漏失、评分错误、结构冲突和依据不足能够分开诊断。分析结果是否有用，最终还要由具体应用回答。下一章将把这些语言线索用于实体、关系和事件的事实抽取，并进一步要求输出受任务本体与原文证据约束；没有表达或没有足够背景支持的信息，不能因为某种分析形式允许填写，就被当成事实补入。
